\documentclass[ppgc,pep,oneside,noabntcite]{infufrgs}

% -----------------------------------------------------------------------------
% Pacotes adicionais
% -----------------------------------------------------------------------------
\usepackage[utf8]{inputenc}
\usepackage{graphicx}
\usepackage{amsmath}
\usepackage{url}
\usepackage{enumitem}
\usepackage{lmodern}

\setlist[itemize]{
label=\textbullet,
leftmargin=1.5em,
itemsep=2pt,
parsep=0pt
}

\setlist[enumerate]{
leftmargin=1.7em,
itemsep=2pt,
parsep=0pt
}

% -----------------------------------------------------------------------------
% Metadados
% -----------------------------------------------------------------------------
\title{Aprendizado por Reforço Multiobjetivo Condicionado por Preferências para Seleção Online de Métodos HTN em NPCs de Jogos}
\translatedtitle{Preference-Conditioned Multi-Objective Reinforcement Learning for Online HTN Method Selection in Game NPCs}

\author{Filho}{Anderson G. A. C.}
\advisor[Prof.~Dr.]{Nunes Alegre}{Lucas}
\date{julho}{2026}

\keyword{Planejamento hierárquico}
\keyword{HTN}
\keyword{Aprendizado por reforço multiobjetivo}
\keyword{Políticas condicionadas por preferências}
\keyword{Game AI}
\keyword{NPCs adaptativos}

\translatedkeyword{Hierarchical planning}
\translatedkeyword{HTN}
\translatedkeyword{Multi-objective reinforcement learning}
\translatedkeyword{Preference-conditioned policies}
\translatedkeyword{Game AI}
\translatedkeyword{Adaptive NPCs}

\begin{document}

\maketitle

\begin{abstract}
Este Plano de Estudos e Pesquisa propõe uma arquitetura híbrida para agentes
não jogáveis na qual o planejamento por Redes Hierárquicas de Tarefas
(
\textit{Hierarchical Task Networks} -- HTN) preserva a validade semântica do
comportamento, enquanto o Aprendizado por Reforço Multiobjetivo
(
\textit{Multi-Objective Reinforcement Learning} -- MORL) adapta a seleção de
métodos de decomposição às prioridades atuais do agente. Diferentemente da
formulação inicial do projeto, o MORL não atua sobre ações motoras ou primitivas.
Sua ação corresponde à escolha, em tempo de execução, de um método HTN entre as
alternativas cuja aplicabilidade já foi validada simbolicamente.

A arquitetura é composta por três elementos. O HTN define tarefas, métodos,
pré-condições, restrições e regras de decomposição. Um Gerador de Pesos
Orientado a Objetivos produz um vetor dinâmico de preferências a partir do
objetivo, estado, personalidade e contexto do agente. Uma política MORL
condicionada por preferências estima o valor multiobjetivo dos métodos
aplicáveis e seleciona o melhor compromisso sob o vetor vigente. As tarefas
primitivas resultantes permanecem sob responsabilidade de controladores
tradicionais, como máquinas de estados, Behavior Trees, navegação ou scripts.

A hipótese central é que essa integração pode superar prioridades HTN estáticas
e heurísticas de utilidade fixas em cenários dinâmicos, mantendo validade
semântica, interpretabilidade, controle de designer e desempenho compatível com
execução em tempo real. A avaliação será conduzida em um ambiente controlado de
jogo, com recompensas vetoriais, mudanças contextuais, preferências dinâmicas,
baselines simbólicos, estudos de ablação e métricas de desempenho da tarefa,
qualidade multiobjetivo, adaptação, diversidade comportamental, validade,
interpretabilidade e custo computacional.
\end{abstract}

\begin{translatedabstract}
This Study and Research Plan proposes a hybrid architecture for non-player
characters in which Hierarchical Task Network (HTN) planning preserves the
semantic validity of behaviour, while Multi-Objective Reinforcement Learning
(MORL) adapts decomposition-method selection to the agent's current
priorities. Unlike the initial formulation of the project, MORL does not act on
motor or primitive actions. Its action is the online selection of an HTN
method among alternatives whose applicability has already been validated by
the symbolic model.

The architecture consists of three components. HTN defines tasks, methods,
preconditions, constraints, and decomposition rules. A Goal-Oriented Weight
Generator produces a dynamic preference vector from the agent's goal, state,
personality, and context. A preference-conditioned MORL policy estimates the
multi-objective value of applicable methods and selects the best trade-off
under the current vector. The resulting primitive tasks remain under the
control of traditional game systems, such as finite-state machines, Behavior
Trees, navigation, or scripts.

The central hypothesis is that this integration can outperform static HTN
priorities and fixed utility heuristics in dynamic scenarios while preserving
semantic validity, interpretability, designer control, and real-time decision
performance. The evaluation will use a controlled game-like environment with
vector rewards, contextual changes, dynamic preferences, symbolic baselines,
ablation studies, and metrics covering task performance, multi-objective
quality, adaptability, behavioural diversity, validity, interpretability, and
computational cost.
\end{translatedabstract}

\tableofcontents

% =============================================================================
\chapter{Identificação e Delimitação do Plano}
% =============================================================================

Este documento apresenta o Plano de Estudos e Pesquisa vinculado ao mestrado
em Ciência da Computação no PPGC/UFRGS. A proposta investiga a integração entre
planejamento HTN e aprendizado por reforço multiobjetivo para a tomada de
decisão de NPCs em jogos digitais.

\section{Dados gerais}

\begin{description}
\item[Aluno:] Anderson G. A. C. Filho.
\item[Programa:] Mestrado em Ciência da Computação, PPGC/UFRGS.
\item[Orientador:] Prof. Dr. Lucas Nunes Alegre.
\item[Projeto institucional:] \textit{Efficient Reinforcement Learning}.
\item[Tema:] seleção online de métodos de decomposição HTN por MORL condicionado por preferências.
\item[Domínio de aplicação:] tomada de decisão hierárquica de NPCs em ambientes de jogo controlados.
\item[Estado inicial:] baseline HTN funcional integrado a um GridWorld com execução por ticks, sensores, navegação e replanejamento.
\end{description}

\section{Contexto de pesquisa}

NPCs são frequentemente controlados por técnicas simbólicas e baseadas em
regras, como máquinas de estados, Behavior Trees \cite{colledanchise2018}, GOAP
e HTN, amplamente discutidas em Game AI \cite{millington2019}. Essas técnicas
fornecem previsibilidade, consistência semântica e controle ao projetista, mas
normalmente dependem de prioridades fixas ou heurísticas construídas
manualmente. A hibridização entre GOAP e HTN também aparece como alternativa
para planejamento em narrativas interativas \cite{artar2019}.
Em ambientes dinâmicos, o número de condições necessárias para representar todas
as variações de contexto pode crescer rapidamente.

O aprendizado por reforço permite adaptação por experiência, porém o
aprendizado direto de ações de baixo nível amplia os espaços de estado e ação,
pode exigir grandes volumes de interação e reduz a transparência da decisão.
O MORL acrescenta a possibilidade de representar explicitamente objetivos
conflitantes, como sobrevivência, cumprimento da missão, proteção de aliados,
conservação de recursos, furtividade e velocidade.

A proposta combina essas propriedades: o HTN define o que é semanticamente
válido; uma fonte de preferências fornece $\mathbf{w}_t$ para expressar o
compromisso vigente; e o MORL escolhe qual método HTN válido deve decompor a
tarefa composta corrente.
A geração de preferências é comparada como componente auxiliar; codificadores
são implementações opcionais dessa fonte, e não uma camada de planejamento
independente.

\section{Problema de pesquisa}

Planejadores HTN tradicionais selecionam métodos por ordem fixa, custos
manuais ou heurísticas estáticas.
Embora previsíveis, esses mecanismos podem ser inadequados quando a melhor
decomposição depende de saúde, munição, condição de aliados, força inimiga,
perigos, objetivo da missão, personalidade ou eventos ocorridos durante a execução.

Uma política puramente aprendida, por outro lado, pode selecionar
comportamentos inválidos, incoerentes ou difíceis de explicar.
O problema investigado é, portanto, como introduzir adaptação multiobjetivo na
escolha de métodos HTN sem perder a validade semântica, a rastreabilidade e a auditabilidade
e o controle de designer fornecidos pelo planejamento simbólico.

\section{Questão principal}

\begin{quote}
O aprendizado por reforço multiobjetivo condicionado por preferências pode
melhorar a seleção online de métodos de decomposição HTN para NPCs, preservando
validade semântica, rastreabilidade, auditabilidade e desempenho de decisão em tempo real?
\end{quote}

\section{Questões secundárias}

\begin{enumerate}
\item A arquitetura produz melhores compromissos multiobjetivo que prioridades HTN estáticas?
\item Uma única política pode adaptar decisões a diferentes vetores $\mathbf{w}_t$ sem retreinamento?
\item Um grafo relacional ou um modelo profundo melhora a inferência de $\mathbf{w}_t$ em relação a regras explícitas?
\item A filtragem simbólica reduz a complexidade de aprendizado e impede métodos inválidos?
\item As decisões permanecem consistentes com as prioridades correntes do agente?
\item Qual é o custo computacional da consulta MORL durante a decomposição online?
\end{enumerate}

\section{Delimitação do escopo}

O foco está no nível hierárquico de decisão.
A ação da política MORL é um método de decomposição HTN, não um comando motor.
Movimento, mira, animação, desvio de obstáculos, pathfinding e execução individual
de ataques serão realizados por controladores determinísticos ou sistemas
tradicionais de jogos.

O escopo experimental inicial será restrito a um agente principal, um domínio
controlado, três a cinco tarefas compostas, duas a cinco alternativas por
tarefa e três a cinco objetivos.
Coordenação multiagente fica fora da contribuição principal.
Uma fonte relacional de preferências é a primeira alternativa estrutural; um
codificador profundo é condicional aos resultados, aos dados e ao orçamento
experimental disponíveis.

% =============================================================================
\chapter{Fundamentação Teórica Inicial}
% =============================================================================

\section{Aprendizado por reforço}

Aprendizado por reforço é uma área de aprendizado de máquina voltada à tomada
de decisão sequencial.
Um agente interage com um ambiente, observa estados, executa ações e recebe
recompensas.
O objetivo é aprender uma política que maximize algum critério de retorno esperado
ao longo do tempo.

Um problema de aprendizado por reforço pode ser descrito pela interação entre
agente e ambiente:

\begin{equation}
s_t \rightarrow a_t \rightarrow r_{t+1} \rightarrow s_{t+1}.
\end{equation}

Nessa interação, $s_t$ é o estado observado no instante $t$, $a_t$ é a ação
escolhida pelo agente, $r_{t+1}$ é a recompensa recebida após a execução da
ação e $s_{t+1}$ é o próximo estado do ambiente.
O agente busca aprender como escolher ações que maximizem recompensas futuras,
considerando que uma decisão tomada no presente pode afetar estados e recompensas
posteriores.

Essa formulação é relevante para a pesquisa porque o projeto investiga agentes
autônomos capazes de tomar decisões em ambientes simulados.
Antes de integrar aprendizado diretamente ao agente, é necessário definir
claramente como estados, ações, recompensas, objetivos e transições serão
representados.
Por isso, a base teórica em Processos de Decisão de Markov, funções de valor,
diferença temporal e aprendizado por reforço multiobjetivo é importante para
as etapas futuras do trabalho.

\section{Processos de Decisão de Markov}

Um Processo de Decisão de Markov, ou MDP, é uma formulação matemática para
problemas de tomada de decisão sequencial sob incerteza.
Essa formulação fornece a base formal para grande parte dos algoritmos de
aprendizado por reforço.

Um MDP pode ser definido como uma tupla:

\begin{equation}
M = \langle S, A, P, R, \gamma \rangle.
\end{equation}

Nessa definição:

\begin{itemize}
\item $S$ é o conjunto de estados possíveis do ambiente;
\item $A$ é o conjunto de ações disponíveis para o agente;
\item $P(s' \mid s,a)$ é a função de transição, que define a probabilidade de ir para o estado $s'$ após executar a ação $a$ no estado $s$;
\item $R(s,a,s')$ é a função de recompensa, que define o retorno imediato associado à transição de $s$ para $s'$ por meio da ação $a$;
\item $\gamma$ é o fator de desconto, com $0 \leq \gamma \leq 1$, usado para controlar a importância de recompensas futuras.
\end{itemize}

A propriedade de Markov estabelece que o próximo estado depende apenas do
estado atual e da ação executada, e não de todo o histórico anterior.
Em outras palavras, se o estado atual contém toda a informação necessária para
prever a dinâmica futura do ambiente, então o histórico completo pode ser ignorado
para fins de decisão.

Formalmente, essa propriedade pode ser escrita como:

\begin{equation}
P(s_{t+1} \mid s_t,a_t,s_{t-1},a_{t-1},\dots,s_0,a_0)
=
P(s_{t+1} \mid s_t,a_t).
\end{equation}

Em um ambiente discreto inicial, como um GridWorld, essa propriedade significa
que o estado precisa conter todas as informações relevantes para a decisão do
agente.
Por exemplo, a posição do agente, a posição de recursos, a posição de objetivos,
o conjunto de obstáculos e variáveis booleanas relevantes devem ser suficientes
para determinar a dinâmica futura do ambiente.

Uma política define o comportamento do agente.
Em uma política estocástica, $\pi(a \mid s)$ representa a probabilidade de
escolher a ação $a$ quando o agente está no estado $s$:

\begin{equation}
\pi(a \mid s).
\end{equation}

Em uma política determinística, a política pode ser representada simplesmente
como uma função que mapeia estados para ações:

\begin{equation}
\pi(s) = a.
\end{equation}

O objetivo do agente é encontrar uma política que maximize o retorno esperado.
O retorno descontado a partir do instante $t$ é definido como:

\begin{equation}
G_t =
R_{t+1}
+
\gamma R_{t+2}
+
\gamma^2 R_{t+3}
+
\dots.
\end{equation}

De forma compacta:

\begin{equation}
G_t =
\sum_{k=0}^{\infty} \gamma^k R_{t+k+1}.
\end{equation}

O fator de desconto $\gamma$ controla a importância de recompensas futuras.
Quando $\gamma$ é próximo de zero, o agente prioriza recompensas imediatas.
Quando $\gamma$ é próximo de um, o agente considera recompensas de longo
prazo.

A função de valor de estado mede o retorno esperado ao seguir uma política
$\pi$ a partir de um estado $s$:

\begin{equation}
V_{\pi}(s) =
E_{\pi}[G_t \mid s_t = s].
\end{equation}

A função de valor de ação mede o retorno esperado ao executar uma ação $a$ em
um estado $s$ e, a partir daí, seguir a política $\pi$:

\begin{equation}
Q_{\pi}(s,a) =
E_{\pi}[G_t \mid s_t = s, a_t = a].
\end{equation}

Para este projeto, os MDPs são relevantes porque fornecem a formulação
matemática da interação entre agente e ambiente.
Eles permitem descrever a dinâmica sequencial em termos de estados, ações,
transições e recompensas.
Em etapas futuras, essa formulação permitirá comparar agentes simbólicos com
agentes baseados em aprendizado.

\section{Q-Learning como base conceitual}

Q-Learning é um algoritmo clássico de aprendizado por reforço usado para
aprender uma função de valor de ação.
Embora não seja o método principal pretendido para esta pesquisa, ele é relevante
como fundamento conceitual, pois introduz ideias centrais que também aparecem em
métodos mais avançados: função de valor, equação de Bellman, diferença temporal e
atualização baseada em experiência.

A função ótima de valor de ação, denotada por $Q^*(s,a)$, representa o retorno
esperado ao executar a ação $a$ no estado $s$ e seguir uma política ótima a
partir do próximo estado.
Essa função satisfaz a equação de otimalidade de Bellman:

\begin{equation}
Q^*(s,a) =
\sum_{s'} P(s' \mid s,a)
\left[
R(s,a,s')
+
\gamma \max_{a'} Q^*(s',a')
\right].
\end{equation}

Essa equação expressa que o valor ótimo de uma ação depende de duas partes:
a recompensa imediata obtida após sua execução e o melhor valor futuro
possível a partir do próximo estado.

Na prática, o agente geralmente não conhece a função de transição
$P(s' \mid s,a)$ nem a função de recompensa completa do ambiente.
Por isso, Q-Learning estima os valores de ação a partir de amostras de experiência.
Uma experiência pode ser representada por:

\begin{equation}
(s_t, a_t, r_{t+1}, s_{t+1}).
\end{equation}

A regra de atualização do Q-Learning é:

\begin{equation}
Q(s_t,a_t)
\leftarrow
Q(s_t,a_t)
+
\alpha
\left[
r_{t+1}
+
\gamma \max_{a'} Q(s_{t+1},a')
-
Q(s_t,a_t)
\right].
\end{equation}

Nessa equação, $\alpha$ é a taxa de aprendizado, $\gamma$ é o fator de
desconto, $r_{t+1}$ é a recompensa observada e
$\max_{a'} Q(s_{t+1},a')$ representa a melhor estimativa de valor futuro no
estado seguinte.

O termo entre colchetes pode ser interpretado como erro de diferença temporal:

\begin{equation}
\delta_t =
r_{t+1}
+
\gamma \max_{a'} Q(s_{t+1},a')
-
Q(s_t,a_t).
\end{equation}

Assim, a atualização também pode ser escrita como:

\begin{equation}
Q(s_t,a_t)
\leftarrow
Q(s_t,a_t)
+
\alpha \delta_t.
\end{equation}

A intuição é que o agente ajusta sua estimativa atual na direção de um alvo
temporal.
Esse alvo é composto pela recompensa imediata mais a melhor estimativa de
valor futuro:

\begin{equation}
\text{target} =
r_{t+1}
+
\gamma \max_{a'} Q(s_{t+1},a').
\end{equation}

Q-Learning é considerado um método \textit{off-policy}.
Isso significa que ele pode aprender o valor da política ótima mesmo
quando as ações são escolhidas por uma política exploratória.
Por exemplo, o agente pode escolher ações usando uma política
$\varepsilon$-\textit{greedy}, mas atualizar sua função $Q$ usando o melhor
valor futuro possível, dado por $\max_{a'} Q(s_{t+1},a')$.

Para este plano, Q-Learning não deve ser tratado como técnica
principal da pesquisa nem como componente previsto da arquitetura
final.
Sua função é servir como base conceitual para entender como
funções de valor podem ser estimadas a partir da interação
com o ambiente.

A técnica de interesse para a continuidade do projeto é aprendizado por
reforço multiobjetivo.
Nesse caso, em vez de aprender uma única função de valor escalar, o
agente deve lidar com múltiplos critérios de decisão, representados por
diferentes sinais de recompensa ou diferentes funções de valor.
Assim, Q-Learning aparece neste plano apenas como ponto de partida
conceitual para compreender funções de valor e atualização temporal,
enquanto MORL representa a direção metodológica pretendida para a camada
de aprendizado.

\section{Multi-armed bandits}

Uma base importante estudada no início da pesquisa foi o problema dos
\textit{multi-armed bandits}.
Nesse problema, o agente escolhe repetidamente entre $k$ ações, cada uma
com uma recompensa de média desconhecida.
O valor real de uma ação é:

\begin{equation}
q_*(a) = E[R_t \mid A_t = a].
\end{equation}

Como esse valor é desconhecido, o agente mantém uma estimativa $Q_t(a)$.
A estimativa por média amostral é:

\begin{equation}
Q_t(a) =
\frac{\sum_{i=1}^{t-1} R_i I(A_i = a)}
{\sum_{i=1}^{t-1} I(A_i = a)}.
\end{equation}

Essa formulação é importante porque introduz conceitos fundamentais como
estimativa de valor, exploração, aproveitamento e atualização incremental em
um cenário mais simples do que MDPs completos.
Enquanto bandits não modelam transições explícitas entre estados, MDPs
incorporam o efeito das ações sobre estados futuros.
Por isso, bandits são uma porta de entrada conceitual para o
estudo de aprendizado por reforço.

\section{Atualização incremental}

A atualização incremental evita recalcular a média completa a cada nova
observação.
A regra geral é:

\begin{equation}
Q_{n+1}(a) = Q_n(a) + \alpha_n(R_n - Q_n(a)).
\end{equation}

O termo $R_n - Q_n(a)$ representa o erro de previsão.
Quando $\alpha_n = 1/n$, a atualização equivale à média amostral.
Quando $\alpha$ é constante, o método dá maior peso a recompensas recentes,
o que é útil em ambientes não estacionários.

Essa ideia reaparece no Q-Learning.
A diferença é que, no caso de Q-Learning, o erro de atualização inclui
não apenas a recompensa imediata, mas também uma estimativa do melhor
valor futuro possível.

\section{Exploração e aproveitamento}

O dilema exploração versus aproveitamento é central em aprendizado por
reforço.
Explorar significa testar ações pouco conhecidas.
Aproveitar significa escolher a ação com melhor valor estimado.

Uma política simples para equilibrar os dois comportamentos é
$\varepsilon$-\textit{greedy}.
Com probabilidade $1-\varepsilon$, o agente escolhe a melhor ação conhecida.
Com probabilidade $\varepsilon$, escolhe uma ação aleatória.

Outra estratégia estudada foi UCB, que escolhe ações considerando tanto o
valor estimado quanto a incerteza:

\begin{equation}
A_t =
\arg\max_a
\left[
Q_t(a) + c\sqrt{\frac{\ln t}{N_t(a)}}
\right].
\end{equation}

Aqui, $N_t(a)$ indica quantas vezes a ação $a$ foi escolhida até o instante
$t$, e $c$ controla a intensidade da exploração.

No projeto, exploração e aproveitamento serão relevantes principalmente nas
etapas futuras de aprendizado.
A camada simbólica atual não aprende por tentativa e erro; ela usa
decomposição explícita de objetivos e ações.
Entretanto, quando métodos de aprendizado forem introduzidos, será necessário
definir como o agente explora ações sem violar restrições simbólicas
importantes.

\section{Planejamento simbólico}

Planejamento simbólico representa estados, tarefas, ações e restrições por
estruturas discretas e interpretáveis.
Em vez de aprender diretamente todo o comportamento, o agente raciocina sobre
pré-condições, efeitos e relações de decomposição.
Essa representação permite incorporar conhecimento de domínio e impedir
decisões incompatíveis com regras narrativas ou de gameplay.

No presente trabalho, as restrições simbólicas são tratadas como
\textit{hard constraints}.
Preferências como segurança, rapidez ou economia de recursos são tratadas
separadamente como critérios flexíveis.
Essa separação é essencial: uma preferência pode mudar em tempo de execução,
enquanto uma regra como não utilizar um recurso indisponível deve permanecer
obrigatória.

\section{Planejamento HTN}

HTN, ou \textit{Hierarchical Task Network Planning}, decompõe tarefas
compostas em subtarefas até produzir tarefas primitivas executáveis
\cite{erol1994,georgievski2014,nau2003shop2}.
Um domínio HTN define tarefas, métodos de decomposição, pré-condições
e restrições.
Para uma tarefa composta $\tau_t$ e estado $s_t$, o conjunto de métodos
aplicáveis é:

\begin{equation}
M_{\mathrm{valid}}(s_t,\tau_t)
=
\left\{m \in M(\tau_t) \mid \mathrm{pre}(m,s_t)=\mathrm{true}\right\}.
\end{equation}

Na arquitetura proposta, o HTN não escolhe imediatamente o primeiro método
aplicável.
Ele produz o conjunto de alternativas válidas e o entrega ao componente
aprendido.
Assim, o modelo simbólico mantém a responsabilidade pela validade, enquanto
a preferência entre alternativas válidas pode ser adaptada por experiência.

\section{Codificação de preferências}

GOAP representa objetivos, ações, pré-condições, efeitos e custos para buscar
sequências que transformem o estado corrente em um estado desejado
\cite{orkin2005,orkin2006}.
Neste trabalho, GOAP não será introduzido como um segundo planejador entre
HTN e MORL. Suas relações entre objetivos, condições, efeitos e recursos
podem, contudo, inspirar a representação de contexto usada por um codificador
de preferências.

O codificador transforma objetivo, estado, perfil e contexto em um vetor
normalizado de preferências:

\begin{equation}
\mathbf{w}_t = E(g_t,s_t,p,c_t),
\qquad
w_{i,t}\geq 0,
\qquad
\sum_i w_{i,t}=1.
\end{equation}

Nessa expressão, $E$ é o codificador, $g_t$ é o objetivo corrente, $s_t$ é o
estado, $p$ representa um perfil comportamental, $c_t$ é o contexto e $w_{i,t}$
é o peso do objetivo $i$ no instante $t$.
Perfis fixos e regras explícitas constituem baselines reproduzíveis.
A principal alternativa estrutural é um grafo relacional inspirado em GOAP,
sem executar sua busca de planejamento; um modelo de aprendizado profundo será
avaliado apenas se os dados, o orçamento e os resultados dos baselines justificarem
essa extensão.
Assim, a comparação entre codificadores é uma parte auxiliar da pesquisa, enquanto
a contribuição central permanece a seleção MORL de métodos HTN válidos.

\section{Aprendizado por reforço multiobjetivo}

No MORL, a recompensa é um vetor
\cite{roijers2013,hayes2022morl}:

\begin{equation}
\mathbf{r}_t =
\left[r_{1,t},r_{2,t},\dots,r_{n,t}\right].
\end{equation}

Cada componente $r_{i,t}$ representa um objetivo, como sobrevivência,
progresso da missão, proteção de aliados, conservação de recursos ou tempo.
A política pode ser condicionada por $\mathbf{w}_t$, permitindo que uma única
aproximação represente diferentes compromissos entre objetivos.
Trabalhos sobre pesos dinâmicos e transferência de políticas fornecem
referências para essa formulação \cite{abels2019,barreto2017successor,alegre2022ols,alegre2026gpi}.

Para cada método válido, a política estima um vetor de valor:

\begin{equation}
\mathbf{Q}(s_t,\tau_t,m,\mathbf{w}_t)
=
\left[Q_1,Q_2,\dots,Q_n\right].
\end{equation}

Aqui, $\tau_t$ é a tarefa composta corrente, $m$ é um método candidato e cada
$Q_i$ estima o retorno descontado do objetivo $i$ ao selecionar $m$. A
dependência explícita de $\mathbf{w}_t$ permite à aproximação aprender valores
ou políticas condicionadas pela preferência, sem transferir ao HTN a decisão de
compromisso entre objetivos.

Sob escalarização linear, a seleção é:

\begin{equation}
 m_t^* =
 \arg\max_{m\in M_{\mathrm{valid}}(s_t,\tau_t)}
 \mathbf{w}_t^{\top}
 \mathbf{Q}(s_t,\tau_t,m,\mathbf{w}_t).
\end{equation}

A máscara de métodos válidos restringe exploração e inferência às alternativas
permitidas pelo HTN. Desse modo, o MORL aprende preferências entre decisões
semanticamente aceitáveis, em vez de aprender diretamente ações de movimento,
mira ou animação.

\section{Decisão hierárquica e abstração temporal}

A seleção de um método HTN pode produzir uma sequência de subtarefas cuja
execução dura vários ticks.
Consequentemente, a decisão aprendida ocorre em uma escala temporal superior à
de ações primitivas e pode ser modelada de forma compatível com abstração
temporal e formulações semi-Markovianas \cite{sutton1999options,lin2024hipo}.
O retorno associado a uma escolha de método deve considerar os resultados obtidos
durante a execução da decomposição, até a próxima decisão hierárquica ou até uma
condição de interrupção.
Se o método selecionado dura $k$ passos primitivos, a unidade de atribuição de
crédito é:

\begin{equation}
\mathbf{R}^{\mathrm{method}}_t =
\sum_{i=0}^{k-1}\gamma^i\mathbf{r}_{t+i}.
\end{equation}

Nessa expressão, $\mathbf{r}_{t+i}$ é a recompensa vetorial no passo primitivo
$t+i$, $\gamma$ é o fator de desconto e $k$ é a duração da execução até a
próxima escolha de método, término ou interrupção.

\section{Arquitetura híbrida proposta}

A integração é definida pelo seguinte fluxo:

\begin{align}
(s_t,\tau_t)
&\longrightarrow \text{filtragem HTN}
\longrightarrow M_{\mathrm{valid}}(s_t,\tau_t),\\
(g_t,s_t,p,c_t)
&\longrightarrow E
\longrightarrow \mathbf{w}_t,\\
(M_{\mathrm{valid}},s_t,\tau_t,\mathbf{w}_t)
&\longrightarrow \text{MORL}
\longrightarrow m_t^*.
\end{align}

O método selecionado é decomposto em subtarefas, e as tarefas primitivas são
executadas por controladores tradicionais.
Após a execução, o ambiente produz um novo estado e um vetor de recompensas,
permitindo atualizar ou avaliar a política.
Essa divisão torna explícitas quatro responsabilidades: validade semântica,
preferência comportamental, utilidade aprendida e execução de baixo nível.

\begin{figure}[htbp]
\centering
\resizebox{0.94\textwidth}{!}{%
\begin{tikzpicture}[
    font=\small,
    >=Latex,
    node distance=8mm and 16mm,
    block/.style={
        rectangle, rounded corners=2mm, draw, thick, align=center,
        minimum height=10mm, minimum width=33mm
    },
    game/.style={block, draw=gamegreen, fill=gamegreen!8},
    htn/.style={block, draw=htnblue, fill=htnblue!8},
    morl/.style={block, draw=morlred, fill=morlred!8},
    runtime/.style={block, draw=envgray, fill=envgray!7},
    arrow/.style={->, thick}
]
    \node[game] (director) {Game / AI Director};
    \node[game, below left=10mm and 19mm of director] (task)
        {Main Task\\$T_{\mathrm{main}} = C_0$};
    \node[game, below right=10mm and 19mm of director] (weight)
        {Preference Vector\\$\mathbf{w}_t=[w_1,\ldots,w_d]$};
    \node[runtime, below=24mm of director] (state)
        {World State\\$WS_t$};
    \node[runtime, left=45mm of state] (sensors)
        {Environment / Sensors};

    \draw[arrow, gamegreen] (director) -- (task);
    \draw[arrow, gamegreen] (director) -- (weight);
    \draw[arrow, envgray, dashed] (sensors) -- (state);

    \node[htn, below=11mm of state] (compound)
        {Current Compound Task\\$C_t$};
    \draw[arrow, htnblue] (task) |- (compound);
    \draw[arrow, htnblue] (state) -- (compound);

    \node[htn, below=10mm of compound] (filter)
        {HTN Symbolic Filter\\Preconditions / Constraints};
    \node[htn, below=9mm of filter] (mask)
        {Feasible-Method Mask\\$\mathcal{M}_{\mathrm{valid}}(C_t,WS_t)$};
    \draw[arrow, htnblue] (compound) -- (filter);
    \draw[arrow, htnblue] (filter) -- (mask);

    \node[morl, below=10mm of mask] (morl)
        {Preference-Conditioned MORL\\$\mathbf{Q}_{\theta}(C_t,WS_t,M_i,\mathbf{w}_t)$};
    \draw[arrow, morlred] (mask) -- (morl);
    \draw[arrow, morlred, dashed] (weight.east)
        to[out=0,in=15] node[pos=.58,right] {preferences} (morl.east);

    \node[morl, below=10mm of morl] (selected)
        {Selected Method\\$M_t^*$};
    \draw[arrow, morlred] (morl) --
        node[right] {$\arg\max U_{\mathbf{w}_t}$} (selected);

    \node[htn, below=10mm of selected] (decompose)
        {HTN Decomposition\\Decompose only $M_t^*$};
    \node[htn, below=10mm of decompose] (network)
        {Resulting Task Network};
    \draw[arrow, htnblue] (selected) -- (decompose);
    \draw[arrow, htnblue] (decompose) -- (network);

    \node[htn, below left=11mm and 19mm of network] (primitive)
        {Primitive Tasks};
    \node[htn, below right=11mm and 19mm of network] (subcompound)
        {Nested Compound Task};
    \draw[arrow, htnblue] (network) -- (primitive);
    \draw[arrow, htnblue] (network) -- (subcompound);
    \draw[arrow, htnblue, dashed, rounded corners] (subcompound.east)
        -- ++(18mm,0) |- node[pos=.25,right,align=left]
        {next HTN\\decision point} (compound.east);

    \node[runtime, below=11mm of primitive] (plan)
        {Final Primitive Plan};
    \draw[arrow, envgray] (primitive) -- (plan);
\end{tikzpicture}%
}
\caption{Arquitetura proposta para seleção online de métodos HTN guiada por MORL.
O HTN determina a máscara de métodos aplicáveis; o MORL seleciona diretamente um
único método $M_t^*$ condicionado por $WS_t$ e $\mathbf{w}_t$. Apenas o método
selecionado é decomposto, e a seleção se repete em subtarefas compostas.}
\label{fig:htn-morl-architecture}
\end{figure}


% =============================================================================
\chapter{Hipóteses e Arquitetura Proposta}
% =============================================================================

\section{Hipótese principal}

A hipótese principal é que uma política MORL condicionada por preferências
pode selecionar métodos de decomposição HTN de forma mais adaptativa que
prioridades estáticas ou heurísticas de pesos fixos, mantendo a validade
semântica garantida pelo HTN.

\section{Hipóteses secundárias}

\begin{description}
\item[H1 --- desempenho adaptativo:] a arquitetura completa obterá melhor
desempenho multiobjetivo sob mudanças ambientais que um HTN com ordem fixa de
métodos.

\item[H2 --- generalização de preferências:] uma política condicionada por
vetores de preferência alterará suas escolhas de maneira consistente quando a
importância relativa dos objetivos mudar, inclusive para combinações não
observadas diretamente no treinamento.

\item[H3 --- validade semântica:] como o HTN filtra os métodos antes da decisão
MORL, a arquitetura completa não selecionará métodos que violem pré-condições
ou restrições definidas pelo domínio.

\item[H4 --- interpretabilidade:] cada escolha poderá ser explicada pela tarefa
composta atual, pelos métodos aplicáveis e rejeitados, pelo vetor de
preferências, pelos valores multiobjetivo estimados e pela utilidade
escalarizada final.

\item[H5 --- redução do espaço de decisão:] restringir o MORL aos métodos HTN
aplicáveis produzirá um espaço de ação menor e semanticamente mais significativo
que o aprendizado direto de ações primitivas.
\end{description}

\section{Componentes da arquitetura}

\subsection{HTN: validade e estrutura}

O HTN define tarefas compostas e primitivas, métodos, pré-condições, requisitos
de recursos e restrições narrativas ou de gameplay. Para cada decisão, ele
produz $M_{\mathrm{valid}}(s_t,\tau_t)$. Métodos impossíveis ou proibidos são
eliminados antes da consulta ao modelo aprendido.

\subsection{Gerador de Pesos Orientado a Objetivos}

O GOWG produz $\mathbf{w}_t$ a partir do objetivo, estado, personalidade e
contexto. Exemplos de componentes são sobrevivência, missão, proteção de
aliados, conservação de recursos, furtividade e velocidade. Regras como
``aumentar o peso de sobrevivência quando a saúde diminui'' serão explicitadas
e versionadas para garantir reprodutibilidade.

\subsection{MORL: seleção de método}

O MORL recebe estado, tarefa composta, vetor de preferência e máscara de
métodos válidos. Sua ação é escolher um método HTN. Portanto, a formulação
corrige a ambiguidade da hipótese inicial: o MORL não é a camada motora e não
substitui os controladores de execução. Ele decide \emph{como decompor} a tarefa
corrente entre alternativas semanticamente permitidas.

\section{Exemplo de decisão}

Para a tarefa composta \texttt{HandleThreat}, o domínio pode definir os métodos
\texttt{Direct\-Engagement}, \texttt{Seek\-Cover},
\texttt{Protect\-Ally}, \texttt{Retreat} e \texttt{Flank}. Se não
houver rota lateral, \texttt{Flank} é removido. Se não houver aliado,
\texttt{Protect\-Ally} é removido. O MORL compara apenas os métodos restantes sob
$\mathbf{w}_t$.

Uma explicação de decisão poderá assumir a forma: o agente selecionou
\texttt{Protect\-Ally} porque o método era aplicável e apresentou a maior
utilidade estimada sob um vetor no qual proteção do aliado e sobrevivência
possuíam os maiores pesos.

\section{Ciclo online}

\begin{enumerate}
\item o HTN identifica a tarefa composta corrente;
\item o HTN calcula os métodos aplicáveis;
\item o GOWG gera o vetor de preferências;
\item o MORL seleciona um método válido;
\item o método é decomposto em subtarefas;
\item controladores tradicionais executam as tarefas primitivas;
\item o ambiente fornece novo estado e recompensa vetorial;
\item o ciclo reinicia na próxima decisão hierárquica.
\end{enumerate}


% =============================================================================
\chapter{Baseline Inicial Baseado em HTN}
% =============================================================================

\section{Motivação para iniciar pelo HTN}

A primeira etapa técnica do plano é construir um baseline HTN standalone. Essa
decisão foi tomada porque o HTN fornece a estrutura de alto nível necessária
para organizar o comportamento do agente.

Antes de integrar uma fonte de preferências e o seletor MORL de métodos, é necessário
garantir a representação de estados simbólicos, a decomposição de tarefas, a
verificação de pré-condições, a aplicação de efeitos, a geração e a execução de
planos, a atualização do estado por sensores e o replanejamento quando o plano
deixa de ser válido.

\section{Componentes implementados}

A implementação atual contém uma classe de estado de mundo e representações de
pré-condições, efeitos, tarefas primitivas, tarefas compostas e métodos de
decomposição. Sobre essas abstrações, o domínio HTN e o \textit{planner} recursivo
produzem planos, enquanto uma interface abstrata de ação e seus estados de
execução permitem ao agente executor realizar as tarefas. O sistema também
inclui sensores, um mecanismo de notificação de mudança de estado, uma abstração
de busca de caminhos com BFS para o grid, o ambiente GridWorld e um
visualizador em terminal.

\section{Estado simbólico}

O estado simbólico representa fatos observáveis sobre o mundo. Ele funciona
como uma estrutura de pares chave-valor. No GridWorld atual, inclui
\texttt{agent\_x}, \texttt{agent\_y}, \texttt{has\_key},
\texttt{door\_open}, \texttt{at\_key}, \texttt{at\_door},
\texttt{at\_goal} e \texttt{done}.

O estado pode ser copiado.
Isso é essencial porque o planner precisa simular efeitos sem modificar o
ambiente real.

\section{Tarefas primitivas}

Uma tarefa primitiva representa uma unidade planejável que pode ser executada
por uma ação concreta.
Formalmente, uma tarefa primitiva pode ser descrita por:

\begin{equation}
t_p = \langle nome, acao, pre, eff \rangle.
\end{equation}

Aqui, $pre$ representa o conjunto de pré-condições e $eff$ representa o
conjunto de efeitos.

\section{Tarefas compostas}

Uma tarefa composta representa uma intenção de alto nível.
Ela não é executada diretamente.
Em vez disso, é decomposta em subtarefas por meio de métodos.

Formalmente:

\begin{equation}
t_c = \langle nome, metodos \rangle.
\end{equation}

\section{Métodos}

Um método representa uma forma possível de decompor uma tarefa composta.
Ele possui pré-condições e uma lista de subtarefas:

\begin{equation}
m = \langle pre(m), subtasks(m) \rangle.
\end{equation}

Um método só pode ser usado quando suas pré-condições são satisfeitas.

\section{Plano}

Um plano é uma sequência de tarefas primitivas:

\begin{equation}
\pi = \langle a_1,a_2,\dots,a_n \rangle.
\end{equation}

Um plano é válido quando cada ação é aplicável no estado produzido pela
simulação das ações anteriores.

\section{Planner recursivo}

O método \texttt{build\_plan} percorre as tarefas raiz do domínio usando uma
cópia do estado observado. Quando a expansão de uma tarefa tem sucesso, suas
tarefas primitivas são anexadas ao plano e o estado simulado resultante é
propagado para a próxima tarefa raiz, como mostra o Algoritmo
\ref{alg:construir-plano}.

\begin{algorithm}[H]
\caption{Método \texttt{build\_plan} do baseline HTN}
\label{alg:construir-plano}
\begin{algorithmic}[1]
\Require domínio $D$, estado observado $WS_0$
\Ensure plano de tarefas primitivas $P$
\State $P \gets \langle\rangle$; $WS \gets \mathrm{Copia}(WS_0)$
\ForAll{$\tau \in D.\mathrm{tasks}$}
  \State $r \gets \Call{PlanejarRecursivamente}{\langle\rangle, WS, \tau}$
  \If{$r \neq \mathrm{Falha}$}
    \State $(P_\tau, WS) \gets r$
    \State $P \gets P \mathbin{\Vert} P_\tau$
  \EndIf
\EndFor
\State \Return $P$
\end{algorithmic}
\end{algorithm}

\section{Planejamento recursivo e backtracking}

O método \texttt{recursive\_planning} expande uma tarefa sobre um estado
simulado. Para uma tarefa primitiva, ele verifica pré-condições, inclui a tarefa
no plano do ramo e aplica seus efeitos apenas à cópia de estado. Para uma tarefa
composta, obtém os métodos aplicáveis e os tenta na ordem definida no domínio.
Se alguma subtarefa falhar, o método inteiro é descartado e o próximo método é
tentado. Esse é o backtracking em nível de método implementado no baseline.

\begin{algorithm}[H]
\caption{Método \texttt{recursive\_planning} com backtracking por método}
\label{alg:planejar-recursivamente}
\begin{algorithmic}[1]
\Require prefixo $P$, estado simulado $WS$, tarefa $\tau$
\Ensure $(P', WS')$ ou $\mathrm{Falha}$
\If{$\tau$ é primitiva}
  \If{$\neg\mathrm{PreCondicoesSatisfeitas}(\tau, WS)$}
    \State \Return $\mathrm{Falha}$
  \EndIf
  \State $P' \gets \mathrm{Copia}(P)$; $WS' \gets \mathrm{Copia}(WS)$
  \State $P' \gets P' \mathbin{\Vert} \langle\tau\rangle$
  \State $WS' \gets \mathrm{AplicarEfeitos}(\tau, WS')$
  \State \Return $(P', WS')$
\EndIf
\If{$\tau$ é composta}
  \ForAll{$m \in \mathrm{MetodosAplicaveis}(\tau, WS)$, na ordem do domínio}
    \State $P_m \gets \mathrm{Copia}(P)$; $WS_m \gets \mathrm{Copia}(WS)$
    \State $falhou \gets \mathrm{falso}$
    \ForAll{$\tau' \in m.\mathrm{subtasks}$}
      \State $r \gets \Call{PlanejarRecursivamente}{P_m, WS_m, \tau'}$
      \If{$r = \mathrm{Falha}$}
        \State $falhou \gets \mathrm{verdadeiro}$; \State \textbf{break}
      \EndIf
      \State $(P_m, WS_m) \gets r$
    \EndFor
    \If{$\neg falhou$}
      \State \Return $(P_m, WS_m)$
    \EndIf
  \EndFor
\EndIf
\State \Return $\mathrm{Falha}$
\end{algorithmic}
\end{algorithm}

% =============================================================================
\chapter{Ambiente Experimental GridWorld}
% =============================================================================

\section{Motivação}

O GridWorld foi escolhido como ambiente inicial por ser simples, discreto,
controlável e adequado para validar planejamento simbólico.
Nesta etapa, o objetivo não é treinar uma política de RL, mas validar a
arquitetura do baseline HTN.

\section{Modelo formal do ambiente}

O ambiente é modelado como um sistema de transição determinístico:

\begin{equation}
\mathcal{E} = \langle S,A,T,R,s_0,G \rangle.
\end{equation}

Nessa formulação, $S$ é o conjunto de estados possíveis, $A$ é o conjunto de
ações discretas, $T$ é a função de transição, $R$ é a recompensa, $s_0$ é o
estado inicial e $G$ é a condição de objetivo. Na versão atual, $R=1$ somente
quando o objetivo é atingido e $R=0$ nos demais passos. O ambiente não fornece
recompensa vetorial nesta etapa; seu papel é validar o ciclo HTN em um cenário
controlado.

\section{Estado do ambiente}

O estado concreto do GridWorld inclui:

\begin{equation}
s =
\langle
p_a,
p_k,
p_d,
p_g,
O,
hasKey,
doorOpen,
done
\rangle.
\end{equation}

Nessa representação, $p_a$, $p_k$, $p_d$ e $p_g$ são, respectivamente, as
posições do agente, da chave, da porta e do objetivo; $O$ é o conjunto de
obstáculos; e $hasKey$, $doorOpen$ e $done$ registram a posse da chave, a
abertura da porta e o término do episódio. A observação retornada pelo ambiente
preserva esses fatos concretos; o sensor os traduz para os predicados usados
pelo domínio HTN.

\section{Ações do ambiente}

O espaço de ações é discreto, com seis operações: mover para cima, direita,
baixo ou esquerda, pegar a chave e abrir a porta. Movimentos que saem do grid,
atingem um obstáculo ou tentam entrar no objetivo com a porta fechada não
alteram a posição. A coleta da chave só é efetiva sobre sua célula, e a abertura
da porta exige que o agente esteja em sua célula e já possua a chave.

O Algoritmo \ref{alg:gridworld-step} descreve o método \texttt{step}, que
processa uma ação, atualiza o término do episódio e retorna o contrato de cinco
valores do Gymnasium.

\begin{algorithm}[H]
\caption{Método \texttt{step} do ambiente GridWorld}
\label{alg:gridworld-step}
\begin{algorithmic}[1]
\Require ação discreta $a$
\Ensure observação $o$, recompensa $r$, término $terminated$, truncamento $truncated$, informação $info$
\If{$done$}
  \State \Return $(\mathrm{Obs}(), 0, \mathrm{verdadeiro}, \mathrm{falso}, \{\})$
\EndIf
\If{$a \in \{\mathrm{cima},\mathrm{direita},\mathrm{baixo},\mathrm{esquerda}\}$}
  \State \Call{TentarMover}{$a$}
\ElsIf{$a = \mathrm{pegarChave}$}
  \State \Call{TentarPegarChave}{}
\ElsIf{$a = \mathrm{abrirPorta}$}
  \State \Call{TentarAbrirPorta}{}
\Else
  \State sinalizar ação inválida
\EndIf
\State $done \gets (p_a=p_g) \land doorOpen$
\State $r \gets 1$ se $done$; caso contrário, $r \gets 0$
\State \Return $(\mathrm{Obs}(), r, done, \mathrm{falso}, \{\})$
\end{algorithmic}
\end{algorithm}

\section{Objetivo do ambiente}

O objetivo é alcançar a posição final com a porta aberta.
Portanto, o episódio termina quando:

\begin{equation}
p_a = p_g
\quad \text{e} \quad
doorOpen = true.
\end{equation}

\section{Configuração}

\texttt{GridWorldConfig} define largura, altura, posições do agente, chave,
porta e objetivo, obstáculos fixos e aleatórios, além dos estados iniciais de
posse da chave e abertura da porta. Uma posição definida como \texttt{None} é
sorteada no reinício do episódio sem sobrepor entidades especiais ou obstáculos.
Assim, posições fixas permitem testes determinísticos, enquanto uma
\textit{seed} torna cenários aleatórios reproduzíveis.

O Algoritmo \ref{alg:gridworld-reset} descreve o método \texttt{reset}.
Ele resolve o layout de cada episódio, restaura os estados iniciais configurados
e produz a primeira observação.

\begin{algorithm}[H]
\caption{Método \texttt{reset} do ambiente GridWorld}
\label{alg:gridworld-reset}
\begin{algorithmic}[1]
\Require semente opcional $seed$
\Ensure observação inicial $o$ e informação $info$
\State inicializar o gerador aleatório com $seed$
\State resolver posições fixas e aleatórias; adicionar obstáculos aleatórios
\State $p_a \gets p_{\mathrm{inicio}}$
\State $hasKey \gets initialHasKey$; $doorOpen \gets initialDoorOpen$
\State $done \gets (p_a=p_g) \land doorOpen$
\State \Return $(\mathrm{Obs}(), \{\})$
\end{algorithmic}
\end{algorithm}

% =============================================================================
\chapter{Domínio HTN do GridWorld}
% =============================================================================

\section{Especificação do domínio HTN}

O capítulo anterior descreve a dinâmica concreta do GridWorld; este capítulo
especifica como essa dinâmica é abstraída em tarefas e métodos HTN. A tarefa
raiz é \texttt{escape\_grid}. As folhas primitivas são
\texttt{go\_to\_key}, \texttt{pickup\_key}, \texttt{go\_to\_door},
\texttt{open\_door} e \texttt{go\_to\_goal}; elas associam uma ação concreta
a pré-condições e efeitos simbólicos. Em particular, as tarefas de navegação
atualizam \texttt{agent\_x}, \texttt{agent\_y} e os predicados
\texttt{at\_key}, \texttt{at\_door} ou \texttt{at\_goal};
\texttt{pickup\_key} produz \texttt{has\_key=true};
\texttt{open\_door} produz \texttt{door\_open=true}; e
\texttt{go\_to\_goal} produz \texttt{done=true}.

As tarefas compostas \texttt{ensure\_has\_key} e
\texttt{ensure\_door\_open} expressam, respectivamente, a obtenção da chave e
a abertura da porta. Cada uma tem um método vazio quando a condição já está
satisfeita, evitando ações artificiais. A tarefa \texttt{escape\_grid} leva
diretamente ao objetivo quando a porta está aberta; caso contrário, primeiro
garante a chave e a abertura da porta. A Figura
\ref{fig:dominio-htn-gridworld} sintetiza essa construção a partir do layout efetivo
do ambiente.

\begin{figure}[H]
\centering
\resizebox{0.92\textwidth}{!}{%
\begin{tikzpicture}[
    font=\footnotesize,
    >=Latex,
    task/.style={rectangle, rounded corners=2mm, draw=htnblue, thick,
        fill=htnblue!8, align=center, minimum width=32mm, minimum height=7mm},
    primitive/.style={rectangle, rounded corners=2mm, draw=gamegreen, thick,
        fill=gamegreen!8, align=center, minimum width=24mm, minimum height=7mm},
    noop/.style={rectangle, rounded corners=2mm, draw=envgray, thick,
        fill=envgray!7, align=center, minimum width=21mm, minimum height=7mm},
    condition/.style={font=\scriptsize, fill=white, inner sep=1.5pt, align=center},
    arrow/.style={->, thick}
]
    % Decomposição da tarefa raiz.
    \node[task] (escape) at (0,0) {\texttt{escape\_grid}};
    \node[primitive] (direct-goal) at (-5.0,-2.1) {\texttt{go\_to\_goal}};
    \node[task] (key) at (2.4,-2.1) {\texttt{ensure\_has\_key}};
    \node[task] (door) at (2.4,-4.8) {\texttt{ensure\_door\_open}};
    \node[primitive] (final-goal) at (2.4,-7.5) {\texttt{go\_to\_goal}};

    \draw[arrow] (escape) -- node[condition, above left, pos=0.57] {\texttt{door\_open}} (direct-goal);
    \draw[arrow] (escape) -- node[condition, above right, pos=0.57] {$\neg$\texttt{door\_open}} (key);
    \draw[arrow] (key) -- (door);
    \draw[arrow] (door) -- (final-goal);

    % Métodos de ensure_has_key, em uma faixa exclusiva.
    \node[noop] (key-noop) at (-0.2,-3.7) {método vazio};
    \node[primitive] (go-key) at (5.0,-3.7) {\texttt{go\_to\_key}};
    \node[primitive] (pick-key) at (8.3,-3.7) {\texttt{pickup\_key}};
    \draw[arrow] (key) -- node[condition, pos=.55, above left] {\texttt{has\_key}} (key-noop);
    \draw[arrow] (key) -- node[condition, pos=.55, above right] {$\neg$\texttt{has\_key}} (go-key);
    \draw[arrow] (go-key) -- (pick-key);

    % Métodos de ensure_door_open, em uma segunda faixa exclusiva.
    \node[noop] (door-noop) at (-0.2,-6.4) {método vazio};
    \node[primitive] (go-door) at (5.0,-6.4) {\texttt{go\_to\_door}};
    \node[primitive] (open-door) at (8.3,-6.4) {\texttt{open\_door}};
    \draw[arrow] (door) -- node[condition, pos=0.66, above left] {\texttt{door\_open}} (door-noop);
    \draw[arrow] (door) -- node[condition, above right, pos=0.67] {$\neg$\texttt{door\_open} $\land$ \texttt{has\_key}} (go-door);
    \draw[arrow] (go-door) -- (open-door);
\end{tikzpicture}
}
\caption{Decomposições do domínio HTN do GridWorld. Nós azuis são tarefas
compostas, nós verdes são tarefas primitivas e nós cinza representam métodos
vazios para subobjetivos já satisfeitos.}
\label{fig:dominio-htn-gridworld}
\end{figure}


Por exemplo, quando o agente não possui a chave e a porta está fechada, a
decomposição produz o plano
\begin{equation}
\begin{aligned}
\pi = \langle {} &\texttt{go\_to\_key},\texttt{pickup\_key},
\texttt{go\_to\_door},\\
&\texttt{open\_door},\texttt{go\_to\_goal} \rangle.
\end{aligned}
\end{equation}
Se a porta já estiver aberta, o método correspondente reduz o plano à navegação
até o objetivo. Essa redução é consequência das pré-condições e dos métodos
vazios do domínio, e não de uma regra especial no ambiente.

% =============================================================================
\chapter{Execução por Agente}
% =============================================================================

\section{Separação entre planejamento e execução}

Uma decisão arquitetural importante foi separar o planner do agente executor.
O planner decompõe tarefas, verifica pré-condições, aplica efeitos apenas ao
estado simulado e retorna folhas primitivas. O agente, por sua vez, mantém o
plano em execução, solicita um novo plano quando necessário, executa a tarefa
corrente e valida o plano remanescente após mudanças observadas no mundo.

\section{Execução por ticks}

A execução ocorre em ciclos discretos chamados ticks.
Em cada tick, o agente pode reconstruir o plano, executar uma etapa da tarefa
corrente e atualizar o plano pendente. Isso é necessário porque ações de
navegação podem durar vários ticks. O Algoritmo \ref{alg:agent-tick} descreve
o método \texttt{tick} atualmente implementado.

\begin{algorithm}[H]
\caption{Método \texttt{tick} do agente HTN}
\label{alg:agent-tick}
\begin{algorithmic}[1]
\Require mundo de execução $W$
\Ensure resultado do tick $R$
\State $replanned \gets \mathrm{falso}$; $planned \gets \langle\rangle$
\If{\Call{DeveReplanejar}{}}
  \State $plan \gets \Call{build\_plan}{}$; $worldStateChanged \gets \mathrm{falso}$
  \State $replanned \gets \mathrm{verdadeiro}$; $planned \gets \Call{NomesDoPlano}{}$
  \If{$plan=\langle\rangle$}
    \State \Return resultado sem tarefa e mensagem ``nenhum plano válido''
  \EndIf
\EndIf
\State $\tau \gets plan[0]$
\If{$\tau$ não é primitiva}
  \State remover $\tau$ do plano; \Return resultado de tarefa ignorada
\EndIf
\State $status \gets \Call{Executar}{\tau.\mathrm{action},W}$
\If{$status=\mathrm{SUCCESS}$}
  \State remover $\tau$ do plano
\ElsIf{$status=\mathrm{FAILURE}$}
  \State $plan \gets \langle\rangle$
\ElsIf{$status \neq \mathrm{RUNNING}$}
  \State sinalizar estado de ação desconhecido
\EndIf
\State \Return $(\tau.\mathrm{name},status,replanned,planned,\Call{NomesDoPlano}{})$
\end{algorithmic}
\end{algorithm}

\section{Status das ações}

Cada ação retorna \texttt{RUNNING} enquanto permanece em execução,
\texttt{SUCCESS} quando conclui a tarefa ou \texttt{FAILURE} quando não pode
concluí-la. Na navegação, por exemplo, o status permanece \texttt{RUNNING} até
o agente alcançar o destino; a ausência de rota resulta em \texttt{FAILURE}.

\section{Resultado de um tick}

O objeto \texttt{AgentTickResult} registra a tarefa executada, o status
retornado, a ocorrência de replanejamento, os nomes das tarefas do plano
gerado e do plano restante, além de uma mensagem opcional. Essas informações
permitem depurar e analisar o ciclo sem acoplar o agente ao visualizador.

% =============================================================================
\chapter{Sensores e Replanejamento}
% =============================================================================

\section{Motivação dos sensores}

O planner não deve acessar diretamente o ambiente concreto.
Em vez disso, ele deve operar sobre um estado simbólico atualizado por sensores.

Um sensor pode ser representado como uma função:

\begin{equation}
\sigma_i:E \rightarrow O_i.
\end{equation}

Aqui, $E$ representa o ambiente e $O_i$ representa uma observação produzida
pelo sensor.

\section{Sensor do GridWorld}

O \texttt{GridWorldSensor} atualiza no \texttt{WorldState} as coordenadas do
agente, da chave, da porta e do objetivo, os fatos \texttt{has\_key},
\texttt{door\_open} e \texttt{done}, e os predicados derivados
\texttt{at\_key}, \texttt{at\_door} e \texttt{at\_goal}. Dessa forma, o
domínio HTN recebe fatos simbólicos sem depender da representação concreta do
ambiente Gymnasium.

\section{Delegate de mudança de estado}

Depois que os sensores atualizam o estado, o \texttt{SensorSystem} notifica os
assinantes por meio de um \textit{delegate}. O agente copia o estado observado,
atualiza a cópia do planner e marca que houve mudança de mundo, como mostrado
na Figura \ref{fig:sensores-replanejamento}.

\begin{figure}[H]
\centering
\begin{tikzpicture}[
    font=\small,
    >=Latex,
    node distance=7mm and 15mm,
    block/.style={rectangle, rounded corners=2mm, draw, thick, align=center,
        text width=42mm, minimum height=9mm},
    runtime/.style={block, draw=gamegreen, fill=gamegreen!8},
    symbolic/.style={block, draw=htnblue, fill=htnblue!8},
    control/.style={block, draw=fallbackorange, fill=fallbackorange!10},
    decision/.style={diamond, aspect=2.1, draw=black!70, thick, align=center,
        fill=black!3, text width=29mm, inner sep=1.5pt},
    arrow/.style={->, thick}
]
    \node[runtime] (env) {Ambiente GridWorld};
    \node[runtime, below=of env] (sensor) {Sensor atualiza fatos simbólicos};
    \node[symbolic, below=of sensor] (state) {\texttt{WorldState} atualizado};
    \node[control, below=of state] (update) {\textit{Delegate} notifica o agente, que atualiza o planner};
    \node[control, below=of update] (validate) {Validar o plano remanescente no próximo tick};
    \node[decision, below=of validate] (valid) {O plano continua válido?};
    \node[symbolic, below left=of valid] (keep) {Manter execução do plano};
    \node[control, below right=of valid] (replan) {Solicitar novo plano HTN};

    \draw[arrow, gamegreen] (env) -- (sensor);
    \draw[arrow, htnblue] (sensor) -- (state);
    \draw[arrow, fallbackorange] (state) -- (update);
    \draw[arrow, fallbackorange] (update) -- (validate);
    \draw[arrow] (validate) -- (valid);
    \draw[arrow, htnblue] (valid) -- node[above left] {sim} (keep);
    \draw[arrow, fallbackorange] (valid) -- node[above right] {não} (replan);
\end{tikzpicture}
\caption{Atualização sensorial e replanejamento preguiçoso. Uma mudança de
estado não descarta imediatamente o plano: o agente valida o plano remanescente
no tick seguinte e só solicita novo planejamento se ele se tornou inválido.}
\label{fig:sensores-replanejamento}
\end{figure}


\section{Validação preguiçosa}

O agente não descarta o plano imediatamente após qualquer mudança de estado.
Em vez disso, valida se o plano restante ainda pode ser executado.
Essa estratégia evita replanejamento desnecessário.

Dado um plano restante:

\begin{equation}
\pi_r =
\langle a_i,a_{i+1},\dots,a_n \rangle,
\end{equation}

o agente simula as tarefas restantes a partir do estado atual. Se todas as
pré-condições forem satisfeitas, o plano continua válido. O Algoritmo
\ref{alg:validar-plano} explicita essa validação, incluindo os efeitos de cada
tarefa primitiva sobre a cópia de estado.

\begin{algorithm}[H]
\caption{Método \texttt{\_is\_plan\_still\_valid}}
\label{alg:validar-plano}
\begin{algorithmic}[1]
\Require plano restante $\pi_r$, estado observado $WS$
\Ensure \texttt{verdadeiro} se $\pi_r$ permanece executável
\State $WS' \gets \mathrm{Copia}(WS)$
\ForAll{$\tau \in \pi_r$}
  \If{$\tau$ não é primitiva}
    \State \textbf{continue}
  \EndIf
  \If{$\neg\mathrm{PreCondicoesSatisfeitas}(\tau,WS')$}
    \State \Return \texttt{falso}
  \EndIf
  \State $WS' \gets \mathrm{AplicarEfeitos}(\tau,WS')$
\EndFor
\State \Return \texttt{verdadeiro}
\end{algorithmic}
\end{algorithm}

\section{Condição de replanejamento}

O replanejamento ocorre quando não existe plano ou quando uma mudança observada
invalida o plano remanescente. Uma falha de ação limpa o plano e, por isso,
provoca replanejamento no tick seguinte. Formalmente:

\begin{equation}
\mathrm{Replanejar}
\iff
(\pi_r=\langle\rangle)
\lor
(\mathrm{mudancaDeMundo}\land\neg\mathrm{Valido}(\pi_r,WS)).
\end{equation}


Esse mecanismo prepara a arquitetura para ambientes mais dinâmicos, mesmo que
o GridWorld inicial seja determinístico.

% =============================================================================
\chapter{Pathfinding e Navegação}
% =============================================================================

\section{Separação entre intenção e movimento}

O HTN define intenções simbólicas, como ir até a chave, ir até a porta ou ir
até o objetivo.
O pathfinder transforma essas intenções em movimentos concretos.

Essa separação evita acoplamento entre o domínio HTN e o algoritmo de
navegação.

\section{Grafo do GridWorld}

O GridWorld pode ser modelado como um grafo:

\begin{equation}
G_{grid} = \langle V,E \rangle.
\end{equation}

O conjunto $V$ contém células livres do grid.
O conjunto $E$ contém conexões entre células ortogonalmente adjacentes.

Uma célula é válida se está dentro dos limites do grid e não pertence ao
conjunto de obstáculos.

\section{Busca em largura}

A implementação atual usa busca em largura, ou BFS. Como as arestas do grid
possuem custo uniforme, BFS encontra o menor caminho em número de passos.

O Algoritmo \ref{alg:bfs-grid} corresponde ao método
\texttt{GridPathfinder.find\_path}. Ele expande vizinhos ortogonais na ordem
cima, direita, baixo e esquerda, ignora células bloqueadas ou já visitadas e
reconstrói o caminho a partir do mapa de predecessores.

\begin{algorithm}[H]
\caption{Método \texttt{find\_path} por busca em largura}
\label{alg:bfs-grid}
\begin{algorithmic}[1]
\Require início $s$, objetivo $g$, contexto $(largura,altura,bloqueadas)$
\Ensure caminho de $s$ até $g$ ou lista vazia
\If{$s=g$}
  \State \Return $[s]$
\EndIf
\State $fronteira \gets \mathrm{Fila}([s])$; $predecessor[s] \gets \emptyset$
\While{$fronteira$ não está vazia}
  \State $atual \gets \Call{RemoverInicio}{fronteira}$
  \If{$atual=g$}
    \State \textbf{break}
  \EndIf
  \ForAll{$vizinho$ ortogonal válido de $atual$}
    \If{$vizinho \in bloqueadas$ \textbf{ ou } $vizinho \in predecessor$}
      \State \textbf{continue}
    \EndIf
    \State $predecessor[vizinho] \gets atual$
    \State \Call{InserirFim}{fronteira,vizinho}
  \EndFor
\EndWhile
\If{$g \notin predecessor$}
  \State \Return $[\ ]$
\EndIf
\State \Return \Call{ReconstruirCaminho}{$predecessor,g$}
\end{algorithmic}
\end{algorithm}

A ação \texttt{NavigateToPositionAction} consulta esse caminho a cada tick,
executa apenas o próximo movimento por meio de \texttt{env.step} e retorna
\texttt{SUCCESS} ao alcançar o alvo. Se o caminho tiver menos de duas células
antes de alcançar o alvo, a ação retorna \texttt{FAILURE}; caso contrário,
retorna \texttt{RUNNING}. A recomputação do caminho torna a navegação reativa
ao estado corrente do ambiente.

% =============================================================================
\chapter{Evolução para o Experimento Multiobjetivo}
% =============================================================================

\section{Papel do baseline atual}

O GridWorld existente permanece útil como prova de conceito do mecanismo HTN,
da execução por ticks, dos sensores, da navegação e do replanejamento.
Ele não constitui, isoladamente, a avaliação final da hipótese.
Sua função é fornecer uma base testável sobre a qual a interface de seleção de
métodos possa ser introduzida sem misturar falhas do planner com falhas de
aprendizado.
Ele ainda não oferece alternativas genuinamente conflitantes, recompensa
vetorial, nem um protocolo de preferência-condicionada; portanto, não é por si
só um benchmark MORL adequado.

\section{Alteração necessária no planner}

O planner atual percorre métodos aplicáveis segundo uma ordem definida no
domínio e usa backtracking quando uma decomposição falha.
Para a arquitetura proposta, a etapa de seleção deverá ser desacoplada da
enumeração fixa.
No caminho online normal, o seletor MORL recebe a máscara dos métodos aplicáveis,
faz uma inferência e escolhe diretamente um único método $M_t^*$; somente $M_t^*$ é
decomposto.
Dada uma tarefa composta, o planner deverá:

\begin{enumerate}
\item identificar todos os métodos cujas pré-condições são satisfeitas;
\item expor uma representação estável desses métodos ao seletor;
\item aplicar uma máscara de validade durante treinamento e inferência;
\item solicitar ao MORL a escolha direta da alternativa preferida;
\item aplicar \emph{method fallback} somente se a decomposição escolhida falhar;
\item registrar métodos rejeitados, valores estimados e método escolhido.
\end{enumerate}

É importante distinguir inaplicabilidade de falha posterior.
A máscara HTN remove métodos que já violam restrições no estado corrente e esses
 métodos nunca chegam ao MORL. Se uma decomposição inicialmente aplicável não
 produz um plano executável após a expansão de suas subtarefas, o método
 selecionado é mascarado temporariamente e o MORL escolhe novamente entre as
 alternativas restantes.
 Essa recuperação não é uma busca DFS sobre uma lista ordenada: é excepcional e
normalmente o ciclo seleciona e decompõe apenas um método.
Quando o mundo muda durante a execução, o agente deve replanejar a partir do
estado observado, e não fazer backtracking sobre um estado simbólico obsoleto.
A separação entre planejamento inicial e replanejamento durante execução é
consistente com trabalhos recentes de replanning HTN sob mudanças de utilidade
e risco \cite{nabakowski2025}.

\begin{figure}[htbp]
\centering
\resizebox{0.92\textwidth}{!}{%
\begin{tikzpicture}[
    font=\small,
    >=Latex,
    node distance=10mm and 18mm,
    block/.style={
        rectangle, rounded corners=2mm, draw, thick, align=center,
        minimum width=35mm, minimum height=10mm
    },
    htn/.style={block, draw=htnblue, fill=htnblue!8},
    morl/.style={block, draw=morlred, fill=morlred!8},
    fallback/.style={block, draw=fallbackorange, fill=fallbackorange!10},
    runtime/.style={block, draw=envgray, fill=envgray!7},
    decision/.style={
        diamond, aspect=2.1, draw=black!70, fill=black!3, thick, align=center,
        inner sep=1.5pt
    },
    arrow/.style={->, thick}
]
    \node[htn] (compound) {Compound Task $C_t$};
    \node[htn, below=of compound] (filter) {HTN filters methods};
    \node[morl, below=of filter] (select) {MORL selects $M_t^*$};
    \node[htn, below=of select] (decompose) {Decompose only $M_t^*$};
    \node[decision, below=12mm of decompose] (success)
        {Decomposition\\successful?};
    \node[runtime, below left=15mm and 22mm of success] (execute)
        {Execute resulting\\primitive tasks};

    \node[fallback, below right=15mm and 22mm of success] (mask)
        {Mask failed $M_t^*$};
    \node[fallback, below=of mask] (remaining)
        {Updated feasible-method mask};
    \node[decision, below=of remaining] (left)
        {Any methods\\left?};
    \node[runtime, below right=13mm and 18mm of left] (failure)
        {Compound Task\\Failure};

    \node[decision, below=of execute] (invalid)
        {Runtime state\\invalidates plan?};
    \node[runtime, below left=14mm and 19mm of invalid] (continue)
        {Continue execution};
    \node[runtime, below right=14mm and 19mm of invalid] (replan)
        {Sense current $WS$\\and replan};

    \draw[arrow, htnblue] (compound) -- (filter);
    \draw[arrow, morlred] (filter) -- (select);
    \draw[arrow, htnblue] (select) -- (decompose);
    \draw[arrow] (decompose) -- (success);
    \draw[arrow] (success) -- node[above left] {yes} (execute);
    \draw[arrow, fallbackorange] (success) --
        node[above right] {no} (mask);
    \draw[arrow, fallbackorange] (mask) -- (remaining);
    \draw[arrow, fallbackorange] (remaining) -- (left);
    \draw[arrow, fallbackorange, rounded corners] (left.east)
        -- ++(26mm,0) |- node[pos=.22,right,align=left]
        {yes: fallback\\selection} (select.east);
    \draw[arrow] (left) -- node[above right] {no} (failure);

    \draw[arrow] (execute) -- (invalid);
    \draw[arrow] (invalid) -- node[above left] {no} (continue);
    \draw[arrow] (invalid) -- node[above right] {yes} (replan);
    \draw[arrow, envgray, dashed, rounded corners] (replan.east)
        -- ++(28mm,0) |- node[pos=.22,right,align=left]
        {current\\World State} (compound.east);
\end{tikzpicture}%
}
\caption{Separação entre o caminho online normal, a recuperação por
\emph{method fallback} e o replanejamento. O fallback só ocorre após falha de
decomposição e retorna ao MORL com uma máscara atualizada. Mudanças durante a
execução provocam replanejamento a partir do estado corrente, não backtracking
sobre um estado simbólico obsoleto.}
\label{fig:htn-morl-fallback}
\end{figure}


\section{Interface prevista}

A interface mínima entre os componentes deverá conter:

\begin{itemize}
\item estado ou representação de características $\phi(s_t)$;
\item identificador da tarefa composta $\tau_t$;
\item conjunto e máscara de métodos válidos;
\item vetor normalizado de preferências $\mathbf{w}_t$;
\item método selecionado $m_t$;
\item recompensa vetorial acumulada durante a decomposição;
\item duração da decisão e condição de término.
\end{itemize}

\section{Ambiente experimental alvo}

A próxima versão do ambiente deverá representar um cenário tático ou de
sobrevivência controlado.
Exemplos de tarefas compostas incluem defender uma posição, proteger um aliado,
obter acesso a uma área, neutralizar uma ameaça ou recuar sob recursos limitados.
Eventos estocásticos e mudanças de prioridade devem ser configuráveis e
reproduzíveis por sementes.

Uma evolução inspirada em ambientes multiobjetivo como
\textit{Resource Gathering} pode ser utilizada para validar recompensas
vetoriais e ferramentas de treinamento \cite{resourcegathering,mogymnasium,morlbaselines}.
Entretanto, o experimento principal deverá conter escolhas hierárquicas com
múltiplos métodos, pois a contribuição científica está na seleção online de
métodos HTN, e não apenas na navegação multiobjetivo.

Ambientes mais complexos, como Craftax, ficam reservados para validação futura
de escalabilidade e não fazem parte do escopo mínimo do mestrado
\cite{craftax2024,craftaxgithub}. Para padronizar a interação entre ambiente e
algoritmos, a implementação experimental deverá manter uma interface compatível
com Gymnasium \cite{gymnasium}.

\section{Instrumentação requerida}

Cada decisão deverá registrar estado, tarefa composta, métodos aplicáveis,
métodos rejeitados e respectivas pré-condições, vetor de preferências, vetor de
valor estimado por método, utilidade escalarizada, método escolhido, sequência
de subtarefas, recompensa vetorial, duração e resultado.
Esses registros sustentarão tanto a análise quantitativa quanto as explicações
qualitativas.

% =============================================================================
\chapter{Estado Atual da Implementação}
% =============================================================================

\section{Componentes disponíveis}

Até agosto de 2026, a arquitetura completa ainda não está implementada.
O baseline disponível já reúne a representação simbólica do estado,
pré-condições, efeitos, tarefas primitivas e compostas, métodos de
decomposição, domínio HTN e um \textit{planner} recursivo com \textit{backtracking}.
Ele também inclui o agente executor orientado a \textit{ticks}, sensores para
atualização do estado, replanejamento inicial, uma abstração de \textit{pathfinding}
com BFS para o GridWorld, visualizador em terminal e documentação técnica
inicial.
A seleção de métodos aplicáveis agora é uma extensão explícita por
meio de \texttt{MethodSelectionStrategy}: DFS preserva a ordem declarada do
domínio, a estratégia heurística permite ranqueamento definido pela aplicação,
e a estratégia RL fornece um gancho para uma implementação concreta.

\section{Interpretação correta do estágio atual}

O código existente demonstra a infraestrutura simbólica e operacional
necessária para decompor e executar tarefas.
O baseline executável ainda usa DFS nos exemplos e mantém o backtracking do
planner após a ordenação dos métodos.
O agente conserva uma cópia das tarefas-raiz e a entrega ao planner em cada
tentativa de replanejamento.
Não existe, neste estágio, política MORL treinada, vetor dinâmico de
preferências integrado ao planner ou máscara de métodos utilizada por um
modelo aprendido.

Portanto, o baseline deve ser reportado como infraestrutura inicial, e não como
validação da hipótese completa.

\section{Lacunas para a nova formulação}

Para integrar MORL, a estratégia de seleção deverá transformar os métodos
aplicáveis e o estado simbólico em entradas estáveis para a política, mantendo
uma correspondência inequívoca entre as ações da política e os métodos candidatos.
A integração também exige uma máscara de validade, um vetor de preferências
$\mathbf{w}_t$, objetivos e recompensas vetoriais, além da definição do horizonte
de crédito de cada decisão hierárquica.

Perfis fixos, regras explícitas e codificadores relacionais serão avaliados como
alternativas de representação, e não como requisitos da interface do
\textit{planner}. Logs de decisão e testes unitários e de integração darão suporte
à confiabilidade da implementação. Experimentos automatizados, cenários e sementes
repetidas, baselines de ordem fixa, prioridade estática e utilidade manual, além de
um protocolo estatístico, comporão a validação científica da proposta.

\section{Compatibilidade do baseline com a proposta}

A implementação atual não precisa ser descartada.
Tarefas compostas, métodos, pré-condições, efeitos, sensores e replanejamento
são diretamente reutilizáveis.
A principal mudança arquitetural já realizada é transformar a seleção de método
em um ponto de extensão explícito.
A futura entrada de preferência deverá ser integrada a uma estratégia concreta,
sem substituir as pré-condições simbólicas nem o backtracking do HTN.
A execução de tarefas primitivas continua separada e permanece determinística
no escopo inicial.

% =============================================================================
\chapter{Objetivos}
% =============================================================================

\section{Objetivo geral}

Projetar e avaliar uma arquitetura híbrida de Game AI na qual uma política de
aprendizado por reforço multiobjetivo condicionada por preferências seleciona,
em tempo de execução, métodos de decomposição HTN entre as alternativas
aplicáveis.

\section{Objetivos específicos}

\begin{enumerate}
\item definir uma interface formal entre HTN, codificação de preferências e MORL;
\item modelar restrições obrigatórias por pré-condições e regras de aplicabilidade HTN;
\item representar preferências comportamentais flexíveis por vetores dinâmicos de pesos;
\item desenvolver uma política MORL que selecione métodos HTN aplicáveis;
\item implementar um ambiente controlado de jogo com recompensas vetoriais e eventos reproduzíveis;
\item comparar a proposta com mecanismos tradicionais de seleção de métodos;
\item medir desempenho, adaptação, diversidade, validade, interpretabilidade e custo computacional;
\item avaliar generalização para vetores de preferência e condições ambientais não vistos;
\item analisar a contribuição de cada componente por estudos de ablação;
\item comparar perfis fixos, regras explícitas e codificadores relacionais para
inferência de $\mathbf{w}_t$, mantendo fixa a política MORL quando aplicável.
\end{enumerate}

% =============================================================================
\chapter{Metodologia}
% =============================================================================

\section{Revisão da literatura}

Será conduzida uma revisão estruturada sobre HTN, GOAP como fonte de
representações relacionais, MORL, políticas condicionadas por preferências,
aprendizado por reforço hierárquico, arquiteturas neuro-simbólicas, NPCs
adaptativos e Game AI explicável.
A revisão deverá posicionar explicitamente a contribuição no nível de seleção
de métodos HTN e identificar trabalhos que aprendem políticas de baixo nível,
métodos de planejamento ou preferências.

\section{Formalização}

A formalização deverá separar quatro elementos: validade semântica,
preferência comportamental, utilidade aprendida e execução de baixo nível.
Serão definidos o estado, as tarefas, os métodos, a função de aplicabilidade, o
vetor de recompensas, o vetor de preferências, a política condicionada e os
critérios de término de uma decisão hierárquica.

\section{Ambiente e dados}

Os dados serão gerados por simulação.
Cada episódio registrará estados, tarefas, métodos, recompensas, preferências,
resultados, consumo de recursos, dano, tempo e latência.
O ambiente deverá permitir cenários repetíveis, configuração de eventos estocásticos,
execução automatizada, sementes determinísticas e visualização de traços de decisão.

O domínio inicial deverá conter um agente principal, três a cinco tarefas
compostas, duas a cinco decomposições por tarefa e três a cinco objetivos.
Um cenário tático ou de sobrevivência é adequado por produzir compromissos
claros entre missão, segurança, aliados, recursos e tempo.

\section{Implementação HTN}

O HTN definirá a estrutura válida do comportamento.
O planner será adaptado para calcular $M_{\mathrm{valid}}(s_t,\tau_t)$ e permitir que um
seletor externo escolha a alternativa.
As restrições obrigatórias permanecerão exclusivamente na camada simbólica.

\section{Codificação de preferências}

Perfis fixos e regras explícitas converterão estado e contexto em um vetor
normalizado e servirão como baselines.
Exemplos incluem aumentar sobrevivência sob baixa saúde, proteção do aliado
quando seu estado se torna crítico e economia de recursos quando a munição
diminui.
Em seguida, será avaliado um codificador baseado em grafo relacional que
represente objetivos, condições, efeitos, recursos e ameaças; essa representação
é inspirada em GOAP, mas não executa um planejador GOAP. Um modelo profundo será
incluído somente se houver dados, orçamento e evidência experimental de que a
complexidade adicional é justificável.
Todos os codificadores serão avaliados sob a mesma interface $E(g_t,s_t,p,c_t)\mapsto\mathbf{w}_t$.

\section{Treinamento MORL}

A política receberá estado, tarefa composta, máscara de métodos válidos e vetor
de preferência.
A ação será a escolha de um método.
O treinamento amostrará estados iniciais, configurações de cenário, eventos e
vetores de preferência.
O orçamento de treinamento, os hiperparâmetros e as escalas de recompensa serão
controlados e reportados.

A atribuição de crédito será definida no nível da decisão hierárquica.
O retorno poderá acumular recompensas durante a execução das subtarefas produzidas
pelo método até a próxima escolha HTN, término da decomposição ou interrupção por
replanejamento.

\section{Baselines}

\begin{description}
\item[Baseline 1 --- primeiro aplicável:] seleciona o primeiro método válido na ordem do domínio.
\item[Baseline 2 --- prioridade estática:] utiliza prioridades manuais constantes.
\item[Baseline 3 --- utilidade com pesos fixos:] combina avaliações manuais por um único vetor fixo.
\item[Baseline 4 --- utilidade dinâmica baseada em regras:] usa regras explícitas para $\mathbf{w}_t$, mas valores manuais, sem aprendizado.
\item[Baseline 5 --- MORL com perfis ou regras:] usa a política MORL com um codificador simples de preferências.
\end{description}

Esses baselines separam o ganho produzido por preferências dinâmicas do ganho
produzido por valores aprendidos.
Codificadores relacionais e profundos serão comparados à configuração MORL equivalente,
sem alterar a máscara HTN.

\section{Estudos de ablação}

Serão avaliadas, quando tecnicamente viáveis, as seguintes configurações:

\begin{itemize}
\item MORL com preferências fixas;
\item MORL sem filtragem HTN, em ambiente seguro ou com penalidades para invalidez;
\item HTN com perfis fixos ou regras, sem MORL;
\item HTN e MORL sem mudanças dinâmicas de preferência;
\item arquitetura HTN--codificador de preferências--MORL;
\item codificador relacional versus regras explícitas sob a mesma política MORL.
\end{itemize}

\section{Procedimento experimental}

Cada arquitetura será executada sobre conjuntos equivalentes de cenários e
sementes.
Serão mantidos, tanto quanto possível, os mesmos horizontes, orçamentos
computacionais e protocolos de ajuste.
Os logs serão agregados por arquitetura e condição.
A análise incluirá estatística descritiva, intervalos de confiança, testes
adequados à distribuição dos dados e tamanhos de efeito.

Mudanças de contexto e preferência serão introduzidas em pontos controlados do
episódio para medir adaptação.
Vetores não usados diretamente no treinamento e cenários com combinações novas
serão reservados para avaliação de generalização.

% =============================================================================
\chapter{Métricas de Avaliação}
% =============================================================================

\section{Desempenho da tarefa}

\begin{itemize}
\item taxa de sucesso da missão e sobrevivência;
\item conclusão de objetivos e sobrevivência de aliados;
\item dano recebido e recursos consumidos;
\item duração do episódio e número de replanejamentos.
\end{itemize}

\section{Desempenho multiobjetivo}

As métricas serão selecionadas conforme as hipóteses do cenário e reportadas
com suas referências, em particular o protocolo de benchmarking MORL
\cite{felten2023toolkit}.

\begin{itemize}
\item retorno vetorial acumulado;
\item utilidade escalarizada sob diferentes preferências;
\item hipervolume e cobertura de soluções somente quando houver conjunto de soluções
e ponto de referência explicitamente definidos;
\item utilidade esperada sobre uma distribuição de preferências;
\item \textit{Maximum Utility Loss} (MUL), quando houver solução de referência;
\item \textit{Inverted Generational Distance} (IGD), apenas em domínios com
fronteira de Pareto de referência conhecida;
\item desempenho sob vetores de preferência não vistos.
\end{itemize}

Para utilidade linear, a utilidade esperada será reportada apenas quando a
distribuição sobre preferências estiver explicitamente definida; métricas de
hipervolume não substituem a avaliação sob preferências condicionadas.

\section{Adaptação e generalização}

\begin{itemize}
\item desempenho após mudanças contextuais;
\item latência de adaptação após alteração de preferências;
\item frequência e adequação das mudanças de método;
\item recuperação após eventos inesperados;
\item desempenho em cenários e combinações não observados no treinamento.
\end{itemize}

\section{Diversidade e coerência comportamental}

\begin{itemize}
\item distribuição dos métodos HTN selecionados;
\item número de estratégias distintas observadas;
\item variação entre personalidades e vetores de preferência;
\item consistência sob estados e preferências equivalentes;
\item taxa de oscilação ou troca prematura de estratégia.
\end{itemize}

\section{Validade semântica}

\begin{itemize}
\item taxa de seleção de método inválido;
\item violações de pré-condições;
\item tentativas de uso de recursos indisponíveis;
\item violações de sequência lógica de tarefas.
\end{itemize}

Para a arquitetura completa, espera-se taxa de seleção inválida igual a zero,
pois os métodos não aplicáveis são removidos antes da decisão MORL. Falhas de
decomposição, falhas de execução e replanejamentos devem ser contabilizados
separadamente.

\section{Desempenho computacional}

\begin{itemize}
\item latência total de decisão, tempo da fonte de preferências e tempo de inferência;
\item tempo de filtragem HTN;
\item tempo de decomposição do método selecionado;
\item taxa e custo de \emph{method fallback};
\item tempo de treinamento;
\item quantidade de métodos decompostos no caminho normal;
\item consumo de memória;
\item estabilidade da frequência de decisão em tempo real.
\end{itemize}

\section{Densidade de decisão}

Além de medir latência, a avaliação reportará a densidade de decisão: número de
consultas ao seletor por episódio, por segundo simulado e por ação primitiva.
Essa métrica testa se a abstração HTN reduz a frequência efetiva de escolhas sem
presumir, a priori, que seu espaço de decisão é globalmente menor.

\section{Rastreabilidade e auditabilidade}

Cada decisão será registrada com tarefa composta, métodos aplicáveis, métodos
rejeitados e pré-condições falhas, vetor de preferências, valores objetivos,
utilidade escalarizada e método escolhido.
A avaliação verificará completude dos traços e coerência entre prioridades e
decisão, complementada por inspeção qualitativa de episódios representativos.


% =============================================================================
\chapter{Cronograma Preliminar}
% =============================================================================

\section{Fase 1 --- fundamentação e posicionamento}

\begin{itemize}
\item concluir revisão estruturada de HTN, MORL e arquiteturas híbridas;
\item consolidar lacuna, questão de pesquisa e hipóteses;
\item revisar a especificação do domínio experimental.
\end{itemize}

\section{Fase 2 --- formalização}

\begin{itemize}
\item formalizar estado, tarefas, métodos, preferências e recompensas;
\item definir a unidade temporal de cada decisão MORL;
\item especificar as interfaces e os registros de explicação.
\end{itemize}

\section{Fase 3 --- estabilização do baseline}

\begin{itemize}
\item concluir testes do HTN, sensores, execução e replanejamento;
\item desacoplar seleção de método da ordem fixa;
\item implementar máscara e identificadores estáveis de métodos.
\end{itemize}

\section{Fase 4 --- ambiente multiobjetivo}

\begin{itemize}
\item implementar o cenário tático ou de sobrevivência;
\item definir eventos, recompensas vetoriais e sementes;
\item implementar instrumentação e execução automatizada.
\end{itemize}

\section{Fase 5 --- preferências}

\begin{itemize}
\item implementar o GOWG determinístico;
\item validar normalização e mudanças contextuais;
\item criar perfis e condições de teste.
\end{itemize}

\section{Fase 6 --- MORL e integração}

\begin{itemize}
\item selecionar e implementar o algoritmo MORL;
\item integrar máscara de métodos válidos;
\item treinar a política sob múltiplas preferências;
\item validar o ciclo online HTN--GOWG--MORL.
\end{itemize}

\section{Fase 7 --- baselines e ablações}

\begin{itemize}
\item implementar os quatro baselines;
\item implementar configurações de ablação;
\item validar equivalência de cenários e orçamento.
\end{itemize}

\section{Fase 8 --- avaliação}

\begin{itemize}
\item executar experimentos com múltiplas sementes;
\item testar mudanças de preferência e generalização;
\item analisar desempenho, validade, diversidade e custo;
\item produzir tabelas, curvas e traços de decisão.
\end{itemize}

\section{Fase 9 --- dissertação e publicação}

\begin{itemize}
\item consolidar metodologia e resultados;
\item discutir limitações e ameaças à validade;
\item redigir a dissertação e um artigo científico;
\item preparar artefatos de reprodutibilidade.
\end{itemize}

% =============================================================================
\chapter{Riscos e Estratégias de Mitigação}
% =============================================================================

\section{Escopo excessivo}

Integrar HTN, geração de preferências, MORL, simulação rica e coordenação
multiagente pode exceder o escopo do mestrado. A contribuição principal será
restrita à seleção online de métodos HTN para um agente. Multiagente e gerador
de pesos aprendido permanecem opcionais.

\section{Recompensas esparsas ou atrasadas}

Uma escolha hierárquica pode afetar resultados somente após várias tarefas
primitivas. Serão avaliadas recompensas intermediárias por objetivo, janelas de
acumulação e estratégias de \textit{reward shaping}, sempre documentando o
efeito sobre a política.

\section{Aprendizado instável sob preferências variadas}

A política pode não generalizar em todo o simplex de pesos. A mitigação inclui
começar com poucos objetivos, normalizar escalas, usar amostragem controlada e
comparar desempenho em vetores vistos e não vistos.

\section{Oscilação comportamental}

Mudanças frequentes de preferência podem causar trocas repetidas de método.
Serão considerados períodos mínimos de compromisso, custo de troca, histerese
ou limiares de decisão. Essas intervenções serão tratadas como parte explícita
do desenho experimental.

\section{Ambiguidade na atribuição de crédito}

A duração variável das decomposições pode dificultar associar retornos à
escolha do método. O protocolo definirá claramente início, término,
interrupção e desconto da decisão hierárquica, com testes em cenários curtos
antes do domínio completo.

\section{Novidade científica insuficiente}

Há trabalhos que combinam planejamento e aprendizado por reforço. A mitigação
é realizar revisão estruturada e posicionar a contribuição de forma precisa na
interseção entre filtragem HTN, preferências dinâmicas e MORL condicionado para
seleção online de métodos.

\section{Comparação experimental injusta}

Diferenças de orçamento, ajuste ou dificuldade podem favorecer uma
arquitetura. Serão usados cenários pareados, mesmas sementes de avaliação,
orçamentos comparáveis, critérios explícitos de ajuste e relato de tamanhos de
efeito.

% =============================================================================
\chapter{Contribuições Esperadas}
% =============================================================================

\begin{itemize}
\item uma arquitetura formal integrando HTN, fontes dinâmicas de preferências e MORL;
\item um método para empregar MORL especificamente na seleção de métodos de decomposição HTN;
\item uma separação explícita entre restrições simbólicas obrigatórias e preferências comportamentais flexíveis;
\item uma política condicionada a ser avaliada quanto à variação de comportamento sem retreinamento por vetor de pesos;
\item um mecanismo rastreável de decisão hierárquica para NPCs adaptativos;
\item uma comparação experimental entre seleção aprendida e heurísticas tradicionais de HTN;
\item uma comparação controlada entre perfis, regras, grafos relacionais e,
quando viável, codificadores profundos de preferências;
\item evidências sobre benefícios, limitações e custo da integração neuro-simbólica proposta;
\item um framework experimental reutilizável para pesquisas futuras em Game AI híbrida.
\end{itemize}

% =============================================================================
\chapter{Considerações Finais}
% =============================================================================

A reformulação apresentada resolve a principal ambiguidade da hipótese
anterior.
O MORL não é responsável por execução motora nem pela escolha direta de ações
primitivas.
Ele atua no ponto em que um planejador HTN possui mais de um método semanticamente
aplicável e precisa escolher qual decomposição melhor atende às prioridades correntes.

A arquitetura pode ser resumida da seguinte forma: o HTN determina o que é
válido; uma fonte de preferências fornece $\mathbf{w}_t$ para expressar o
compromisso corrente; e o MORL seleciona o método válido mais apropriado.
A geração de preferências é um componente auxiliar e comparável: perfis e
regras são baselines; um grafo relacional inspirado em GOAP e um modelo profundo
são implementações opcionais da fonte, condicionadas ao orçamento e aos resultados.
A execução das tarefas primitivas permanece sob responsabilidade de sistemas
convencionais.

O baseline HTN existente é compatível com esse direcionamento, mas deverá ser
adaptado para expor o conjunto de métodos aplicáveis e aceitar um seletor
externo.
A avaliação deverá demonstrar não apenas ganho de desempenho, mas também
generalização entre preferências, validade semântica, qualidade das explicações
e custo compatível com decisões em tempo real.

\vspace{3em}
\noindent Porto Alegre, agosto de 2026.

\vspace{4em}
\begin{flushleft}
\rule{7cm}{0.4pt}\\
\textbf{Anderson G. A. C. Filho}\\
Mestrando em Ciência da Computação --- PPGC/UFRGS
\end{flushleft}


% =============================================================================
% Referências
% =============================================================================

\begin{thebibliography}{99}

\bibitem{sutton2018}
Sutton, R. S.; Barto, A. G.
\textit{Reinforcement Learning: An Introduction}.
2. ed. Cambridge: MIT Press, 2018.

\bibitem{georgievski2014}
Georgievski, I.; Aiello, M.
\textit{An Overview of Hierarchical Task Network Planning}.
arXiv:1403.7426, 2014.

\bibitem{erol1994}
Erol, K.; Hendler, J.; Nau, D. S.
\textit{HTN Planning: Complexity and Expressivity}.
In: \textit{Proceedings of the Twelfth National Conference on Artificial
Intelligence}. AAAI Press, 1994.

\bibitem{nau2003shop2}
Nau, D.; Au, T.-C.; Ilghami, O.; Kuter, U.; Murdock, J. W.;
Wu, D.; Yaman, F.
\textit{SHOP2: An HTN Planning System}.
Journal of Artificial Intelligence Research, v. 20,
p. 379--404, 2003.

\bibitem{hogg2010}
Hogg, C.; Kuter, U.; Muñoz-Avila, H.
\textit{Learning Methods to Generate Good Plans: Integrating HTN Learning and
Reinforcement Learning}.
In: \textit{Proceedings of the Twenty-Fourth AAAI Conference on Artificial
Intelligence}, p. 1530--1535, 2010.

\bibitem{ghallab2004}
Ghallab, M.; Nau, D.; Traverso, P.
\textit{Automated Planning: Theory and Practice}.
San Francisco: Elsevier, 2004.

\bibitem{artar2019}
Artar, E.
\textit{Goal-Oriented Hierarchical Task Networks and Its Application on
Interactive Narrative Planning}.
Master's thesis, Sabanci University, 2019.

\bibitem{orkin2005}
Orkin, J.
\textit{Agent Architecture Considerations for Real-Time Planning in Games}.
In: \textit{Proceedings of the First AAAI Conference on Artificial
Intelligence and Interactive Digital Entertainment}, 2005.

\bibitem{orkin2006}
Orkin, J.
\textit{Three States and a Plan: The AI of F.E.A.R.}
In: \textit{Game Developers Conference}, 2006.

\bibitem{colledanchise2018}
Colledanchise, M.; Ögren, P.
\textit{Behavior Trees in Robotics and AI}.
Boca Raton: CRC Press, 2018.

\bibitem{kuric2024}
Kuric, D.; Infante, G.; Gomez, V.; Jonsson, A.; van Hoof, H.
\textit{Planning with a Learned Policy Basis to Optimally Solve Complex Tasks}.
In: \textit{Proceedings of the 34th International Conference on Automated
Planning and Scheduling}, 2024.

\bibitem{prabhakar2025}
Prabhakar, N.; Singh, R.; Kokel, H.; Natarajan, S.; Tadepalli, P.
\textit{Combining Planning and Reinforcement Learning for Solving Relational
Multiagent Domains}.
In: \textit{Proceedings of the 24th International Conference on Autonomous
Agents and Multiagent Systems}, 2025.

\bibitem{millington2019}
Millington, I.; Funge, J.
\textit{Artificial Intelligence for Games}.
3. ed. Boca Raton: CRC Press, 2019.

\bibitem{hayes2022morl}
Hayes, C. F.; Rădulescu, R.; Bargiacchi, E.; Källström, J.; Macfarlane, M.;
Reymond, M.; Verstraeten, T.; Zintgraf, L. M.; Dazeley, R.; Heintz, F.;
Howley, E.; Irissappane, A. A.; Mannion, P.; Nowé, A.; Ramos, G.; Restelli, M.;
Vamplew, P.; Roijers, D. M.
\textit{A Practical Guide to Multi-Objective Reinforcement Learning and
Planning}.
Autonomous Agents and Multi-Agent Systems, v. 36, n. 1, 2022.
DOI: \url{https://doi.org/10.1007/s10458-022-09552-y}.

\bibitem{roijers2013}
Roijers, D. M.; Vamplew, P.; Whiteson, S.; Dazeley, R.
\textit{A Survey of Multi-Objective Sequential Decision-Making}.
Journal of Artificial Intelligence Research, v. 48,
p. 67--113, 2013.


\bibitem{abels2019}
Abels, A.; Roijers, D. M.; Lenaerts, T.; Nowé, A.; Steckelmacher, D.
\textit{Dynamic Weights in Multi-Objective Deep Reinforcement Learning}.
In: \textit{Proceedings of the 36th International Conference on Machine
Learning}, PMLR, v. 97, p. 11--20, 2019.

\bibitem{alegre2026gpi}
Alegre, L. N.; Bazzan, A. L. C.; Roijers, D. M.; Nowé, A.; da Silva, B. C.
\textit{Generalized Policy Improvement for Efficient and Robust
Multi-Objective Reinforcement Learning}.
Autonomous Agents and Multi-Agent Systems, v. 40, n. 1, 2026.
DOI: \url{https://doi.org/10.1007/s10458-026-09736-w}.

\bibitem{felten2023toolkit}
Felten, F.; Alegre, L. N.; Nowé, A.; Bazzan, A. L. C.; Talbi, E.-G.;
Danoy, G.; da Silva, B. C.
\textit{A Toolkit for Reliable Benchmarking and Research in
Multi-Objective Reinforcement Learning}.
In: \textit{Proceedings of the 37th Conference on Neural Information
Processing Systems}, 2023.

\bibitem{li2024mohtn}
Li, M.; Liu, X.; Jiang, G.; Liu, W.
\textit{A Novel Hierarchical Task Network Planning Approach for Multi-Objective
Optimization}.
Expert Systems with Applications, v. 251, art. 124058, 2024.
DOI: \url{https://doi.org/10.1016/j.eswa.2024.124058}.

\bibitem{bahrami2025}
Bahrami, S.
\textit{Enhancing HTN Planning with Deep Reinforcement Learning for Method
Selection}.
Master's thesis, University of Stuttgart, 2025.

\bibitem{nabakowski2025}
Nabakowski, L.
\textit{Replanning in Risk Aware HTN Planning}.
Master's thesis, University of Stuttgart, 2025.

\bibitem{mogymnasium}
Farama Foundation.
\textit{MO-Gymnasium: Multi-Objective Reinforcement Learning Environments}.
Disponível em: \url{https://mo-gymnasium.farama.org/}.
Acesso em: julho de 2026.

\bibitem{morlbaselines}
Alegre, L. N.
\textit{MORL-Baselines: Multi-Objective Reinforcement Learning Algorithms}.
Disponível em: \url{https://github.com/LucasAlegre/morl-baselines}.
Acesso em: julho de 2026.

\bibitem{resourcegathering}
Farama Foundation.
\textit{Resource Gathering Environment}.
Disponível em:
\url{https://mo-gymnasium.farama.org/environments/resource-gathering/}.
Acesso em: julho de 2026.

\bibitem{barreto2017successor}
Barreto, A.; Dabney, W.; Munos, R.; Hunt, J. J.; Schaul, T.;
van Hasselt, H.; Silver, D.
\textit{Successor Features for Transfer in Reinforcement Learning}.
In: \textit{Advances in Neural Information Processing Systems},
v. 30, 2017.

\bibitem{barreto2020fast}
Barreto, A.; Hou, S.; Borsa, D.; Silver, D.; Precup, D.
\textit{Fast Reinforcement Learning with Generalized Policy Updates}.
Proceedings of the National Academy of Sciences,
v. 117, n. 48, p. 30079--30087, 2020.

\bibitem{borsa2019usfa}
Borsa, D.; Barreto, A.; Quan, J.; Mankowitz, D. J.;
van Hasselt, H.; Munos, R.; Silver, D.; Schaul, T.
\textit{Universal Successor Features Approximators}.
In: \textit{Proceedings of the 7th International Conference on Learning
Representations}, 2019.

\bibitem{alegre2022ols}
Alegre, L. N.; Bazzan, A. L. C.; da Silva, B. C.
\textit{Optimistic Linear Support and Successor Features as a Basis for
Optimal Policy Transfer}.
In: \textit{Proceedings of the 39th International Conference on Machine
Learning}, p. 394--413, 2022.

\bibitem{sutton1999options}
Sutton, R. S.; Precup, D.; Singh, S.
\textit{Between MDPs and Semi-MDPs: A Framework for Temporal Abstraction in
Reinforcement Learning}.
Artificial Intelligence, v. 112, n. 1--2,
p. 181--211, 1999.

\bibitem{lin2024hipo}
Lin, Y.-A.; Lee, C.-T.; Yang, C.-H.; Liu, G.-T.; Sun, S.-H.
\textit{Hierarchical Programmatic Option Framework}.
In: \textit{Advances in Neural Information Processing Systems}, v. 37,
p. 126677--126724, 2024.
DOI: \url{https://doi.org/10.52202/079017-4024}.

\bibitem{barreto2019optionkeyboard}
Barreto, A.; Borsa, D.; Hou, S.; Comanici, G.; Aygün, E.;
Hamel, P.; Hunt, J. J.; Mourad, S.; Silver, D.; Precup, D.
\textit{The Option Keyboard: Combining Skills in Reinforcement Learning}.
In: \textit{Advances in Neural Information Processing Systems},
v. 32, 2019.

\bibitem{alegre2025optionkeyboard}
Alegre, L. N.; Bazzan, A. L. C.; Barreto, A.; da Silva, B. C.
\textit{Constructing an Optimal Behavior Basis for the Option Keyboard}.
In: \textit{Proceedings of the Thirty-ninth Annual Conference on Neural
Information Processing Systems}, 2025.

\bibitem{klissarov2025hrl}
Klissarov, M.; Bagaria, A.; Luo, Z.; Konidaris, G.; Precup, D.;
Machado, M. C.
\textit{Discovering Temporal Structure: An Overview of Hierarchical
Reinforcement Learning}.
arXiv:2506.14045, 2025.
DOI: \url{https://doi.org/10.48550/arXiv.2506.14045}.

\bibitem{vamplew2011}
Vamplew, P.; Dazeley, R.; Berry, A.; Issabekov, R.; Dekker, E.
\textit{Empirical Evaluation Methods for Multiobjective Reinforcement
Learning Algorithms}.
Machine Learning, v. 84, p. 51--80, 2011.

\bibitem{craftax2024}
Matthews, M.; Beukman, M.; Ellis, B.; Samvelyan, M.; Jackson, M.; Coward, S.;
Foerster, J.
\textit{Craftax: A Lightning-Fast Benchmark for Open-Ended Reinforcement
Learning}.
arXiv:2402.16801, 2024.

\bibitem{craftaxgithub}
Matthews, M.
\textit{Craftax: Crafter and NetHack in JAX}.
Disponível em: \url{https://github.com/MichaelTMatthews/Craftax}.
Acesso em: julho de 2026.

\bibitem{gymnasium}
Towers, M.; Terry, J. K.; Kwiatkowski, A.; Balis, J. U.;
Colas, C.; Deleu, T.; Goulão, M.; Kallinteris, A.; Krimmel, M.;
Perez-Vicente, R.; Pierré, A.; Schulhoff, S.; Tai, J. J.;
Sharma, A.; Younis, O.
\textit{Gymnasium}.
Disponível em: \url{https://gymnasium.farama.org/}.
Acesso em: julho de 2026.

\end{thebibliography}


\end{document}
